% universal-spectral-matching.tex
% Universal Spectral Matching: Reducing All Comparison to Computer Vision
% Through Oscillatory Representation and GPU-Parallel Interference
% Author: Kundai Farai Sachikonye
% Date: 2026

\documentclass[11pt,twocolumn,a4paper]{article}

% ─── Packages ────────────────────────────────────────────────────────
\usepackage[utf8]{inputenc}
\usepackage[T1]{fontenc}
\usepackage{lmodern}
\usepackage{amsmath,amssymb,amsthm,mathtools}
\usepackage{physics}
\usepackage{bm}
\usepackage{graphicx}
\usepackage{xcolor}
\usepackage{booktabs}
\usepackage{array}
\usepackage{multirow}
\usepackage{enumitem}
\usepackage{float}
\usepackage{caption}
\usepackage{subcaption}
\usepackage[colorlinks=true,linkcolor=blue!60!black,citecolor=blue!60!black,urlcolor=blue!60!black]{hyperref}
\usepackage[margin=1.8cm,top=2.2cm,bottom=2.2cm]{geometry}
\usepackage{natbib}
\usepackage{fancyhdr}
\usepackage{titlesec}
\usepackage{microtype}
\usepackage{algorithm}
\usepackage{algpseudocode}
\usepackage{listings}
\usepackage{tabularx}
\usepackage{tcolorbox}

% ─── Listings Style ─────────────────────────────────────────────────
\lstdefinestyle{glsl}{
  language=C,
  basicstyle=\ttfamily\scriptsize,
  keywordstyle=\color{blue!70!black}\bfseries,
  commentstyle=\color{green!50!black}\itshape,
  stringstyle=\color{red!60!black},
  numbers=left,
  numberstyle=\tiny\color{gray},
  numbersep=4pt,
  frame=single,
  breaklines=true,
  columns=fullflexible
}

\lstdefinestyle{api}{
  basicstyle=\ttfamily\scriptsize,
  keywordstyle=\color{blue!70!black}\bfseries,
  frame=single,
  breaklines=true,
  columns=fullflexible
}

% ─── Theorem Environments ────────────────────────────────────────────
\newtheorem{theorem}{Theorem}[section]
\newtheorem{lemma}[theorem]{Lemma}
\newtheorem{corollary}[theorem]{Corollary}
\newtheorem{proposition}[theorem]{Proposition}
\newtheorem{definition}[theorem]{Definition}
\newtheorem{remark}[theorem]{Remark}
\newtheorem{axiom}{Axiom}
\newtheorem{example}[theorem]{Example}

% ─── Custom Commands ─────────────────────────────────────────────────
\newcommand{\kb}{k_{\mathrm{B}}}
\newcommand{\Ocat}{\hat{O}_{\mathrm{cat}}}
\newcommand{\Ophys}{\hat{O}_{\mathrm{phys}}}
\newcommand{\Sk}{S_k}
\newcommand{\St}{S_t}
\newcommand{\Se}{S_e}
\newcommand{\BigO}{\mathcal{O}}
\newcommand{\R}{\mathbb{R}}
\newcommand{\N}{\mathbb{N}}
\newcommand{\Z}{\mathbb{Z}}
\newcommand{\CC}{\mathbb{C}}
\newcommand{\Hilb}{\mathcal{H}}
\newcommand{\Lag}{\mathcal{L}}
\newcommand{\Ham}{\mathcal{H}}
\newcommand{\Part}{\mathcal{P}}
\newcommand{\Cat}{\mathcal{C}}
\newcommand{\taup}{\tau_{\mathrm{p}}}
\newcommand{\tauem}{\tau_{\mathrm{em}}}
\newcommand{\Spec}{\mathrm{Spec}}
\DeclareMathOperator{\Tr}{Tr}
\DeclareMathOperator{\diag}{diag}
\DeclareMathOperator{\sgn}{sgn}
\DeclareMathOperator{\Vis}{Vis}
\DeclareMathOperator{\Match}{Match}

% ─── Page Style ──────────────────────────────────────────────────────
\pagestyle{fancy}
\fancyhf{}
\fancyhead[L]{\small\textit{Universal Spectral Matching}}
\fancyhead[R]{\small\thepage}
\renewcommand{\headrulewidth}{0.4pt}

% ─── Title Section Formatting ────────────────────────────────────────
\titleformat{\section}{\large\bfseries}{\thesection.}{0.5em}{}
\titleformat{\subsection}{\normalsize\bfseries}{\thesubsection.}{0.5em}{}
\titleformat{\subsubsection}{\normalsize\itshape}{\thesubsubsection.}{0.5em}{}

% ═════════════════════════════════════════════════════════════════════
\begin{document}

\twocolumn[
\begin{@twocolumnfalse}
\begin{center}

{\LARGE\bfseries Universal Spectral Matching:\\[4pt]
Reducing All Comparison to Computer Vision\\[4pt]
Through Oscillatory Representation\\[4pt]
and GPU-Parallel Interference}

\vspace{12pt}

{\large Kundai Farai Sachikonye}

\vspace{6pt}

{\small\itshape AIMe Registry for Artificial Intelligence \\[0.3em]
\texttt{kundai.sachikonye@bitspark.com}}

\vspace{6pt}

{\small March 2026}

\vspace{14pt}

\begin{minipage}{0.92\textwidth}
\textbf{Abstract.}
We prove from first principles that every comparison problem---matching
molecules, signals, structures, sequences, or arbitrary data---reduces
to a computer vision problem. The argument proceeds through a chain of
mathematical identities, not analogies. First, we prove the Oscillatory
Necessity Theorem: every persistent dynamical system confined to bounded
phase space necessarily exhibits oscillatory dynamics, with static,
monotonic, and non-recurrent chaotic modes eliminated by boundedness,
Poincar\'{e} recurrence, and non-degeneracy. Second, we prove via
Koopman operator theory that every such oscillatory system admits a
complete spectral decomposition $\Spec(X) = \{(\omega_k, A_k, \phi_k)\}$
into finitely many frequency--amplitude--phase triples. Third, we prove
the Spectral Image Theorem: this spectrum is isomorphic to a
two-dimensional image $I_X : \R^2 \to \R$, with frequency mapped to
horizontal position, phase to vertical position, and amplitude to pixel
intensity, where the mapping is injective and preserves metric structure.
Fourth, we prove the Universal Reduction Theorem: the categorical
distance between any two bounded systems $X$ and $Y$ equals the $L^2$
image distance between their spectral images, establishing an exact
identity $d(X,Y) = d_{\mathrm{CV}}(I_X, I_Y)$. Fifth, we show that GPU
fragment shaders implement massively parallel per-pixel interference,
and that comparing two spectral images reduces to computing the
visibility of their interference pattern in a single render pass.
We specify a complete five-pass GPU shader pipeline---the Universal
Spectral Matching Module---with domain-specific encoders for microscopy,
molecular spectra, chromatography, time series, text, genomic sequences,
and graphs. The pipeline achieves $\BigO(1)$ wall-clock time per
comparison on commodity GPU hardware. We provide the full architectural
blueprint: API specifications, GLSL shader pseudocode, framebuffer
layouts, and platform notes for WebGL\,2, Vulkan, and Metal, together
with validation frameworks against established benchmarks. This paper
serves as the self-contained theoretical foundation and engineering
blueprint for the Universal Spectral Matching module.

\vspace{8pt}
\noindent\textbf{Keywords:} spectral matching, oscillatory dynamics, bounded phase space,
GPU computing, fragment shaders, interference, computer vision,
partition coordinates, harmonic coupling, Koopman operator, similarity metrics
\end{minipage}

\vspace{18pt}
\end{center}
\end{@twocolumnfalse}
]

% ═════════════════════════════════════════════════════════════════════
%                    PART I: THEORETICAL FOUNDATIONS
% ═════════════════════════════════════════════════════════════════════

\section{Introduction}
\label{sec:introduction}

The comparison of objects is among the most fundamental operations in
computation, science, and cognition. Whether one seeks to determine
whether two molecular spectra originate from the same compound, whether
two image patches depict the same structure, whether two text passages
convey the same meaning, or whether two genomic sequences encode
homologous proteins, the underlying operation is \emph{matching}:
quantifying the degree to which two representations are equivalent.

Despite the universality of this need, the computational landscape for
comparison is fragmented. Molecular informatics deploys Tanimoto
coefficients over fingerprint vectors \citep{willett1998, tanimoto1958}.
Computer vision employs feature descriptors such as SIFT \citep{lowe2004}
and ORB \citep{rublee2011}, or learned representations from deep networks
\citep{oquab2023}. Natural language processing computes cosine similarity
over embedding vectors \citep{vaswani2017, reimers2019}. Sequence
bioinformatics uses alignment-based scoring via BLAST \citep{altschul1990}
and BLAT \citep{kent2002}. Time series analysis relies on dynamic time
warping \citep{sakoe1978, berndt1994}. Graph comparison uses
Weisfeiler--Leman tests or graph neural networks \citep{kipf2017, xu2019}.
Each domain has developed its own metrics, its own representations, and
its own computational pipelines, with no unifying theory explaining why
comparison works at all or whether a single computational primitive could
serve all domains.

This paper establishes such a unifying theory and provides its complete
engineering realisation. The central claim is that \emph{all comparison
reduces to computer vision}. This is not a metaphor. It is a chain of
mathematical identities:

\begin{enumerate}[nosep]
  \item Every bounded persistent system IS oscillations
    (Theorem~\ref{thm:oscillatory_necessity}).
  \item Every oscillatory system has a spectrum: frequency, amplitude,
    phase (Theorem~\ref{thm:koopman_spectral}).
  \item A spectrum IS an image: intensity mapped over a frequency--phase
    domain (Theorem~\ref{thm:spectral_image}).
  \item Comparing two spectra = comparing two images
    (Theorem~\ref{thm:universal_reduction}).
  \item GPU fragment shaders implement massively parallel per-pixel
    interference (Theorem~\ref{thm:gpu_interference}).
  \item Therefore ALL matching reduces to a single GPU render pass.
\end{enumerate}

% ── Figure 1 proposal ──────────────────────────────────────────────
% FIGURE 1: High-level pipeline diagram.
% Left: diverse input data (molecule, image, time series, graph, text,
% DNA sequence). Centre: each mapped to a spectral image in the
% (omega, phi) plane. Right: pairwise interference patterns, with
% constructive = match, destructive = mismatch. GPU chip icon
% surrounding the centre and right stages.

The remainder of the paper is organised as follows.
Part~I (Sections~\ref{sec:axioms}--\ref{sec:commutation}) develops the
complete theoretical foundation from two axioms through the universal
reduction theorem and entropy conservation.
Part~II (Sections~\ref{sec:pixel_wavefunctions}--\ref{sec:complexity})
constructs the interference-based comparison framework and proves its
complexity properties.
Part~III (Sections~\ref{sec:gpu_engine}--\ref{sec:module_api})
provides the complete GPU architecture blueprint.
Part~IV (Sections~\ref{sec:app_bioimage}--\ref{sec:conclusion})
presents applications, validation, and conclusions.


% ═════════════════════════════════════════════════════════════════════
\section{Axioms}
\label{sec:axioms}

We require exactly two axioms, from which the entire framework is derived.

\begin{axiom}[Bounded Phase Space]
\label{ax:bounded}
Every physical or representable system $\mathcal{S}$ that persists for
observation occupies a region $\Omega \subset \R^{2d}$ of phase space
(where $d$ is the number of degrees of freedom) that is bounded:
there exists a finite constant $R > 0$ such that for all states
$\mathbf{x} \in \Omega$, $\|\mathbf{x}\| \leq R$. The region $\Omega$
has finite Liouville measure $\mu(\Omega) < \infty$.
\end{axiom}

\noindent\emph{Justification.} Boundedness is not an arbitrary
restriction but a precondition for observability. Any measurement
apparatus has finite range. Any digital representation has finite bit
depth. Any physical container has finite volume. A system that is
unbounded in phase space either escapes observation (and thus cannot be
compared) or is truncated to a bounded subsystem for practical purposes.
Shannon's channel capacity theorem \citep{shannon1948} requires finite
signal power, which implies bounded amplitude in any signal space.
Heisenberg's uncertainty relation \citep{heisenberg1927} implies that
phase-space cells have minimum volume $(\hbar/2)^d$, so a bounded region
contains finitely many distinguishable states. The finite Liouville
measure condition \citep{liouville1838} ensures that Hamilton's equations
preserve the phase-space volume, which is essential for the recurrence
arguments that follow. In computational terms, any datum stored in finite
memory occupies a bounded region of its representation space.

\begin{axiom}[Categorical Observation]
\label{ax:categorical}
Every act of observation partitions the state space of a system into a
finite number of distinguishable categories. Formally, observation is a
surjective map $\mathcal{O}: \Omega \to \{c_1, c_2, \ldots, c_K\}$
from the (possibly continuous) state space to a finite set of $K$
categorical outcomes, where $K$ is determined by the resolution of the
observing instrument.
\end{axiom}

\noindent\emph{Justification.} This axiom encodes the finite resolution
of all measurement. A CCD camera partitions photon flux into discrete
pixel intensities. A mass spectrometer partitions molecular masses into
discrete $m/z$ bins. A human observer partitions colours into nameable
categories. Von Neumann's measurement postulate \citep{vonneumann1932}
formalises this for quantum systems: measurement projects a continuous
state onto a discrete eigenvalue. Nyquist's sampling theorem
\citep{nyquist1928} establishes that a bandwidth-limited continuous
signal is fully characterised by discrete samples, making discrete
categorisation lossless for bounded systems. The axiom applies equally
to classical instruments and to computational discretisation (floating
point, pixel grids, hash tables).

\begin{remark}
These two axioms are independent. Axiom~\ref{ax:bounded} constrains the
system; Axiom~\ref{ax:categorical} constrains the observer. Together
they define the arena in which all comparison operations take place:
bounded systems observed through finite-resolution instruments.
\end{remark}


% ═════════════════════════════════════════════════════════════════════
\section{The Oscillatory Necessity Theorem}
\label{sec:oscillatory_necessity}

We now prove that oscillation is the unique long-term dynamical mode
available to systems satisfying Axioms~\ref{ax:bounded}
and~\ref{ax:categorical}.

\begin{definition}[Dynamical Modes]
\label{def:modes}
Let $\mathbf{x}(t)$ denote the trajectory of a system in bounded phase
space $\Omega$. We classify the long-term behaviour ($t \to \infty$)
into four exhaustive categories:
\begin{enumerate}[label=(\roman*),nosep]
  \item \textbf{Static}: $\mathbf{x}(t) \to \mathbf{x}^*$ (convergence
    to a fixed point).
  \item \textbf{Monotonic}: $\|\mathbf{x}(t)\|$ is eventually
    monotonically increasing or decreasing.
  \item \textbf{Chaotic (aperiodic, non-recurrent)}: the trajectory
    never returns arbitrarily close to any previously visited state.
  \item \textbf{Oscillatory}: the trajectory is recurrent---for every
    $\varepsilon > 0$ and every visited state $\mathbf{x}(t_0)$, there
    exists $T > 0$ such that
    $\|\mathbf{x}(t_0 + T) - \mathbf{x}(t_0)\| < \varepsilon$.
\end{enumerate}
\end{definition}

\begin{theorem}[Oscillatory Necessity]
\label{thm:oscillatory_necessity}
Let $\mathcal{S}$ be a dynamical system satisfying
Axiom~\ref{ax:bounded} (bounded phase space with finite Liouville
measure) that persists indefinitely ($t \to \infty$) and is
non-degenerate (possesses at least two distinguishable states under
Axiom~\ref{ax:categorical}). Then the long-term dynamics of
$\mathcal{S}$ are oscillatory in the sense of
Definition~\ref{def:modes}(iv).
\end{theorem}

\begin{proof}
We proceed by elimination of the three alternative modes.

\textbf{Elimination of the static mode.}
If $\mathbf{x}(t) \to \mathbf{x}^*$, then after some time $t_c$ the
system permanently occupies a single categorical bin under
Axiom~\ref{ax:categorical}. A system in a single state carries zero
information ($H = 0$ where $H$ is the Shannon entropy
\citep{shannon1948} of the categorical distribution) and cannot be
compared with anything, since comparison requires distinguishing at
least two values. By the non-degeneracy hypothesis, the system possesses
at least two distinguishable states, so the static mode contradicts
persistence of a non-degenerate system. More precisely, a system that
reaches a fixed point ceases to generate new observational data; it has
terminated as a dynamical entity. We require persistence, so the static
mode is excluded.

\textbf{Elimination of the monotonic mode.}
If $\|\mathbf{x}(t)\|$ is eventually monotonically increasing, then
since $\Omega$ is bounded ($\|\mathbf{x}\| \leq R$), the trajectory
must converge to the boundary. By continuity, this convergence is to a
limit point, reducing to the static case on the boundary. If
$\|\mathbf{x}(t)\|$ is eventually monotonically decreasing, the
trajectory converges to a limit point (possibly the origin), again
reducing to the static case. Monotonicity in any phase-space coordinate
is similarly bounded: each coordinate is confined to $[-R, R]$, so a
monotonic coordinate must converge. If all coordinates converge, the
system is static. If some converge and others do not, the
non-convergent coordinates must be non-monotonic, i.e., oscillatory in
at least those dimensions. Thus pure monotonic behaviour is impossible
in bounded phase space; it either terminates (static) or is partial,
leaving oscillatory residual degrees of freedom.

\textbf{Elimination of non-recurrent chaos.}
Consider a trajectory $\mathbf{x}(t)$ in bounded $\Omega$ that never
returns close to any previously visited state. Since $\Omega$ is bounded,
it has finite Lebesgue measure $\mu(\Omega) < \infty$. We invoke the
Poincar\'{e} Recurrence Theorem \citep{poincare1890}: for any
measure-preserving flow on a bounded phase space with finite measure,
almost every trajectory returns arbitrarily close to its initial
condition infinitely often. The finite Liouville measure condition in
Axiom~\ref{ax:bounded} ensures that this theorem applies. If the flow is
not exactly measure-preserving (e.g., weakly dissipative), the trajectory
is attracted to a compact attractor $\mathcal{A} \subset \Omega$ of
lower dimension, on which it is eventually confined. On $\mathcal{A}$,
the dynamics are recurrent by compactness: the sequence
$\{\mathbf{x}(n)\}_{n \in \N}$ in a compact set has a convergent
subsequence (Bolzano--Weierstrass), implying the trajectory returns
arbitrarily close to previously visited states. Thus non-recurrent
trajectories have measure zero in any bounded phase space; for a
persistent system, recurrence is almost certain.

Since non-recurrent chaos has measure zero and is therefore
non-generic, and since static and monotonic modes are excluded by
boundedness and non-degeneracy, the remaining mode---oscillatory
(recurrent) dynamics---is the unique generic long-term behaviour.

To strengthen this: recurrence, by Definition~\ref{def:modes}(iv),
means the system revisits neighbourhoods of previous states, which is
precisely the defining property of oscillation in the generalised sense
(not necessarily periodic, but quasi-periodic or recurrent). The
categorically observed trajectory under Axiom~\ref{ax:categorical}
exhibits a sequence of categorical states that necessarily repeats
patterns, since there are only $K$ categories and the system persists
indefinitely. By the Pigeonhole Principle, every sequence of length
$> K$ on $K$ symbols must contain a repeated symbol, and by recurrence,
these repetitions are structured (not random), constituting oscillatory
patterns.
\end{proof}

\begin{corollary}
\label{cor:spectral_existence}
Every persistent non-degenerate system in bounded phase space admits a
Fourier-type spectral decomposition.
\end{corollary}

\begin{proof}
By Theorem~\ref{thm:oscillatory_necessity}, the system is recurrent.
The categorical time series $c(t) = \mathcal{O}(\mathbf{x}(t))$ is a
bounded function on $\R$ (taking values in $\{1, \ldots, K\}$). By
recurrence, it is almost periodic in the sense of Bohr
\citep{bohr1925}: for every $\varepsilon > 0$, the set of
$\varepsilon$-almost periods has bounded gaps. Bohr's theorem guarantees
that every almost periodic function admits a generalised Fourier series
$c(t) = \sum_k A_k e^{i\omega_k t}$ with countably many frequencies
$\omega_k$ and amplitudes $A_k$. Since the function takes only $K$
values, the spectrum has at most countably many non-zero components, and
by the finite bandwidth implied by Axiom~\ref{ax:bounded}, only
finitely many frequencies carry significant amplitude. Thus the system
admits a finite spectral decomposition.
\end{proof}


% ═════════════════════════════════════════════════════════════════════
\section{Partition Coordinates}
\label{sec:partition_coordinates}

The spectral decomposition of Corollary~\ref{cor:spectral_existence}
provides frequencies and amplitudes but does not specify how the
categorical states are organised hierarchically. We now derive a
coordinate system for spectral components that emerges from the
hierarchical nesting of oscillatory modes within bounded spherical
phase space.

\begin{definition}[Hierarchical Oscillatory Nesting]
\label{def:nesting}
Consider a system with multiple oscillatory modes. The modes are
naturally ordered by frequency: the slowest mode (lowest frequency)
defines the coarsest partition of the trajectory, and successively
faster modes subdivide within each period of the slower mode. We define
the \emph{principal partition number} $n \in \{1, 2, 3, \ldots\}$ as
the index of the frequency tier, ordered from lowest to highest.
\end{definition}

\begin{theorem}[Partition Coordinate System]
\label{thm:partition_coords}
The spectral components of a bounded oscillatory system in a domain with
spherical symmetry are uniquely characterised by four quantum numbers
$(n, \ell, m, s)$:
\begin{enumerate}[label=(\roman*),nosep]
  \item $n \in \{1, 2, 3, \ldots\}$: the principal partition number
    (frequency tier / depth of hierarchical subdivision).
  \item $\ell \in \{0, 1, \ldots, n-1\}$: the angular partition number,
    indexing distinct oscillatory sub-modes within tier $n$ that differ
    in their angular (phase-space rotational) character.
  \item $m \in \{-\ell, -\ell+1, \ldots, \ell\}$: the magnetic
    partition number, indexing the orientation of the angular mode
    relative to an observation axis, giving $2\ell + 1$ projections.
  \item $s \in \{+\tfrac{1}{2}, -\tfrac{1}{2}\}$: the binary partition
    number, encoding the two-fold degeneracy arising from the
    time-reversal symmetry of bounded oscillatory dynamics (a mode and
    its time-reverse are observationally distinguishable by phase).
\end{enumerate}
\end{theorem}

\begin{proof}
\textbf{Derivation of $n$.}
By Axiom~\ref{ax:bounded}, the phase space $\Omega$ is bounded. The
oscillatory decomposition from Corollary~\ref{cor:spectral_existence}
yields a discrete set of frequencies $\omega_1 < \omega_2 < \cdots$.
Since $\Omega$ has finite measure, the number of distinguishable
frequency tiers is bounded (by the Nyquist limit \citep{nyquist1928}
applied to the phase-space resolution from Axiom~\ref{ax:categorical}).
We index these tiers by $n = 1, 2, \ldots, n_{\max}$.

\textbf{Derivation of $\ell$.}
At each frequency tier $n$, the oscillation occurs in a
$(2d)$-dimensional phase space. The angular degrees of freedom of the
oscillatory mode---the directions of oscillation in phase space---are
characterised by the representations of the rotation group $SO(d)$
restricted to the bounded domain. For a $d$-dimensional system, the
angular modes at tier $n$ are indexed by $\ell = 0, 1, \ldots, n-1$.
The upper bound $\ell \leq n - 1$ arises because the $n$th frequency
tier subdivides the domain into $n$ radial zones, and the angular
complexity cannot exceed the radial resolution: an oscillatory mode
with $\ell$ angular nodes requires at least $\ell + 1$ radial tiers
to be resolved. This is the same constraint that arises in the
separation of variables for the Laplacian on a bounded domain
\citep{courant1953}, where the eigenfunction at the $n$th energy level
has angular quantum number bounded by $n - 1$.

\textbf{Derivation of $m$.}
Given angular quantum number $\ell$, the projection onto any fixed
observation axis (chosen by the observer under
Axiom~\ref{ax:categorical}) yields $2\ell + 1$ distinguishable
orientations, indexed by $m = -\ell, -\ell + 1, \ldots, \ell$. This
is the standard result for representations of $SO(3)$, extended to
general $SO(d)$ by the branching rules \citep{weyl1946}. For
$d \geq 3$, additional quantum numbers arise but are suppressed into
the $m$-degeneracy at each $\ell$.

\textbf{Derivation of $s$.}
Time-reversal symmetry $t \to -t$ maps each oscillatory mode to its
reverse. In bounded phase space, the forward and reverse modes are
generically distinguishable (they have opposite phase velocity), giving
a two-fold degeneracy. We encode this as $s = \pm \frac{1}{2}$, a
binary quantum number. This is formally identical to the spin-$\frac{1}{2}$
representation of $SU(2)$, the double cover of $SO(3)$.
\end{proof}

\begin{theorem}[Capacity Formula]
\label{thm:capacity}
The total number of distinguishable spectral states at partition level
$n$ is
\begin{equation}
\label{eq:capacity}
  C(n) = 2n^2.
\end{equation}
\end{theorem}

\begin{proof}
Count all valid $(n, \ell, m, s)$ combinations at level $n$:
\begin{align}
  C(n) &= \sum_{\ell=0}^{n-1} \sum_{m=-\ell}^{\ell} \sum_{s=\pm 1/2} 1
  = 2 \sum_{\ell=0}^{n-1} (2\ell + 1) \nonumber \\
  &= 2 \left[ \sum_{\ell=0}^{n-1} (2\ell + 1) \right]
  = 2 \cdot n^2 = 2n^2,
\end{align}
where we used the identity $\sum_{\ell=0}^{n-1}(2\ell+1) = n^2$.
\end{proof}

\begin{remark}
The capacity $C(n) = 2n^2$ is formally identical to the electron
capacity of atomic shells, but here it arises purely from the
combinatorics of hierarchical oscillatory partitioning in bounded
phase space, without reference to quantum mechanics, electromagnetism,
or atomic structure. The coincidence is explained by the shared
mathematical structure: both atomic orbitals and our partition
coordinates arise from eigenvalue problems on bounded domains with
rotational symmetry \citep{sakurai2017, courant1953}.
\end{remark}


% ═════════════════════════════════════════════════════════════════════
\section{The Triple Equivalence}
\label{sec:triple_equivalence}

We now establish that three superficially different descriptions of a
bounded system---oscillatory dynamics, categorical enumeration, and
statistical partition function---are formally equivalent.

\begin{theorem}[Triple Equivalence]
\label{thm:triple}
For any system $\mathcal{S}$ satisfying Axioms~\ref{ax:bounded}
and~\ref{ax:categorical}, the following three characterisations are
equivalent:
\begin{enumerate}[label=(\alph*),nosep]
  \item \textbf{Oscillatory decomposition}: $\mathcal{S}$ is described
    by spectral components $\{(n_k, \ell_k, m_k, s_k, A_k)\}$
    with frequencies $\omega_k$ and amplitudes $A_k$.
  \item \textbf{Categorical enumeration}: $\mathcal{S}$ is described by
    the occupation numbers $\{N_j\}_{j=1}^{K}$ counting how many
    observations fall into each of $K$ categories, with
    multiplicity $M = \sum_j N_j$.
  \item \textbf{Partition function}: $\mathcal{S}$ is described by the
    partition function $Z = \sum_j e^{-\beta E_j}$ and the associated
    entropy $S = \kb M \ln n$, where $n$ is the effective number of
    occupied partition levels.
\end{enumerate}
Moreover, explicit bijective maps connect all three descriptions.
\end{theorem}

\begin{proof}
\textbf{(a) $\Rightarrow$ (b): Map $\Phi_{ab}$.}
The oscillatory decomposition assigns each spectral component a
partition coordinate $(n, \ell, m, s)$. Under
Axiom~\ref{ax:categorical}, the observer bins these components into
$K$ categories. Define the map $\Phi_{ab}$ by: the occupation number
of category $j$ is
$N_j = \sum_{k: \mathcal{O}(n_k, \ell_k, m_k, s_k) = j} |A_k|^2$
(weighted by spectral power). The total multiplicity is
$M = \sum_j N_j = \sum_k |A_k|^2$ (Parseval's theorem
\citep{fourier1822}). This map is surjective since every category
receives occupation from at least one spectral component (by
non-degeneracy), and injective up to the resolution of the categorical
partition.

\textbf{(b) $\Rightarrow$ (c): Map $\Phi_{bc}$.}
Given occupation numbers $\{N_j\}$, define energies $E_j = -\ln(N_j/M)$
(the surprisal of category $j$) and inverse temperature
$\beta = 1/(\kb T)$ such that $N_j/M = e^{-\beta E_j}/Z$ with
$Z = \sum_j e^{-\beta E_j}$. This is exactly the Boltzmann distribution
\citep{boltzmann1877}. The entropy is
\begin{equation}
\label{eq:entropy}
  S = -\kb \sum_j \frac{N_j}{M} \ln \frac{N_j}{M} = \kb M \ln n,
\end{equation}
where $n = e^{H}$ is the effective number of categories (exponential of
the Shannon entropy $H = -\sum_j p_j \ln p_j$, $p_j = N_j/M$)
\citep{shannon1948}. For uniform occupation, $n = K$ and $S = \kb M \ln K$.
The map $\Phi_{bc}$ is bijective: given $(Z, \beta)$, the occupation
numbers are recovered as $N_j = M \cdot e^{-\beta E_j}/Z$.

\textbf{(c) $\Rightarrow$ (a): Map $\Phi_{ca}$.}
Given the partition function $Z$, each energy level $E_j$ corresponds
to a frequency $\omega_j = E_j / \hbar$ (via the Planck--Einstein
relation, or more generally, by identifying energy with frequency
through the oscillatory dynamics \citep{einstein1905}). The occupation
probability $p_j = e^{-\beta E_j}/Z$ determines the amplitude
$A_j = \sqrt{p_j}$. The phase is determined by the position within the
oscillatory cycle at the time of observation. Thus the partition function
determines a complete oscillatory decomposition. The composition
$\Phi_{ca} \circ \Phi_{bc} \circ \Phi_{ab} = \mathrm{id}$ establishes
the bijective equivalence.
\end{proof}

% ── Figure 2 proposal ──────────────────────────────────────────────
% FIGURE 2: Triple equivalence diagram. Three nodes in a triangle:
% "Oscillatory Decomposition" (top), "Categorical Enumeration"
% (bottom-left), "Partition Function" (bottom-right). Double-headed
% arrows between all three, labelled with the bijective maps
% Phi_ab, Phi_bc, Phi_ca. Centre: "S = k_B M ln n".

\begin{corollary}[Entropy--Multiplicity Relation]
\label{cor:entropy_multiplicity}
The entropy of a bounded system is proportional to the total
multiplicity and the logarithm of the effective partition depth:
\begin{equation}
\label{eq:S_M_n}
  S = \kb \cdot M \cdot \ln(n).
\end{equation}
\end{corollary}


% ═════════════════════════════════════════════════════════════════════
\section{The Fundamental Identity}
\label{sec:fundamental_identity}

\begin{theorem}[Fundamental Identity]
\label{thm:fundamental}
For a bounded oscillatory system observed categorically, the rate of
multiplicity generation equals the oscillation frequency divided by
$2\pi$, which equals the reciprocal of the mean categorical persistence
time:
\begin{equation}
\label{eq:fundamental}
  \frac{dM}{dt} = \frac{\omega}{2\pi} = \frac{1}{\langle\taup\rangle}.
\end{equation}
\end{theorem}

\begin{proof}
Consider the system's categorical trajectory $c(t)$. The multiplicity
$M(t)$ counts the total number of categorical transitions (state changes)
up to time $t$. In one complete oscillation cycle of period
$T = 2\pi/\omega$, the system traverses all $K$ categories and returns
to its starting category (by the recurrence property of
Theorem~\ref{thm:oscillatory_necessity}). Each complete cycle thus
generates one unit of multiplicity (one complete enumeration of the
categorical space). Therefore $dM/dt = 1/T = \omega/(2\pi)$.

The mean categorical persistence time $\langle\taup\rangle$ is the
average duration spent in one complete categorical cycle. Define
$\taup$ as the first return time of the categorical process:
$\taup = \inf\{t > 0 : c(t) = c(0)\}$. By Kac's theorem
\citep{kac1947}, the expected first return time to any state in a
recurrent process equals the reciprocal of the stationary probability
of that state, and the mean return time to \emph{any} previously visited
state is $\langle\taup\rangle = T$ (the period). Thus
$dM/dt = 1/\langle\taup\rangle = \omega/(2\pi)$.
\end{proof}

\begin{remark}
This identity connects three different measurement modalities: counting
(multiplicity $M$), timing (persistence $\taup$), and spectral analysis
(frequency $\omega$). It ensures that any measurement of one quantity
determines the other two, establishing the fundamental interconvertibility
of measurement modalities for bounded systems.
\end{remark}


% ═════════════════════════════════════════════════════════════════════
\section{The Spectral Representation Theorem}
\label{sec:spectral_representation}

We now prove the central theoretical result: every bounded oscillatory
system admits a complete spectral decomposition via Koopman operator
theory.

\begin{definition}[Koopman Operator]
\label{def:koopman}
Let $\Phi^t: \Omega \to \Omega$ denote the flow map of a dynamical
system on bounded phase space $\Omega$. The \emph{Koopman operator}
$\mathcal{U}^t$ is the linear operator acting on the space of
observables $g: \Omega \to \CC$ by composition with the flow:
\begin{equation}
\label{eq:koopman}
  (\mathcal{U}^t g)(\mathbf{x}) = g(\Phi^t(\mathbf{x})).
\end{equation}
\end{definition}

The Koopman operator was introduced by \citet{koopman1931} for
Hamiltonian systems and subsequently extended to general measure-preserving
dynamics. Its key property is that it is \emph{linear}, even when the
underlying dynamics are nonlinear, because it acts on functions of the
state rather than on the state itself \citep{mezic2005, budisic2012}.

\begin{theorem}[Koopman Spectral Decomposition]
\label{thm:koopman_spectral}
Let $\mathcal{S}$ be a persistent non-degenerate system satisfying
Axiom~\ref{ax:bounded}. Then the Koopman operator $\mathcal{U}^t$
associated with $\mathcal{S}$ admits a spectral decomposition: there
exist finitely many eigenvalues $\lambda_k = e^{i\omega_k}$ with
associated eigenfunctions $\varphi_k$ such that for any observable
$g \in L^2(\Omega, \mu)$,
\begin{equation}
\label{eq:koopman_decomp}
  g(\Phi^t(\mathbf{x})) = \sum_{k=1}^{N} A_k \, \varphi_k(\mathbf{x})
    \, e^{i\omega_k t} + r(\mathbf{x}, t),
\end{equation}
where $A_k = \langle g, \varphi_k \rangle$ are the Koopman mode
amplitudes, $\omega_k$ are real frequencies (by the measure-preserving
property), and $r$ is a residual with $\|r\|_{L^2} \to 0$ as
$N \to \infty$.

The spectrum of the system is defined as the finite set
\begin{equation}
\label{eq:spectrum_def}
  \Spec(\mathcal{S}) = \{(\omega_k, A_k, \phi_k)\}_{k=1}^{N},
\end{equation}
where $\phi_k = \arg(A_k)$ is the phase of the $k$th Koopman mode.
\end{theorem}

\begin{proof}
By Axiom~\ref{ax:bounded}, $\Omega$ is bounded with finite measure.
By \citet{koopman1931}, the Koopman operator for a measure-preserving
flow on a space of finite measure is a unitary operator on
$L^2(\Omega, \mu)$. By the spectral theorem for unitary operators on
separable Hilbert spaces \citep{vonneumann1932}, $\mathcal{U}^t$ admits
a spectral decomposition. Since the flow is on a bounded domain,
the eigenvalues lie on the unit circle: $|\lambda_k| = 1$, hence
$\lambda_k = e^{i\omega_k}$ with $\omega_k \in \R$.

The eigenfunction expansion
\eqref{eq:koopman_decomp} follows from the spectral theorem applied to
the unitary group $\{\mathcal{U}^t\}_{t \in \R}$. The amplitudes
$A_k$ are obtained by projecting the observable $g$ onto the
eigenfunctions.

Finiteness of the significant spectral components follows from
Axiom~\ref{ax:categorical}: the observation map $\mathcal{O}$ has
finite range $K$, so the observable $g = \mathcal{O}$ takes only $K$
values. A function with $K$ distinct values requires at most $K$
Fourier modes for exact representation (at each time instant, the
function value is determined by the superposition of these modes).
By the Nyquist criterion \citep{nyquist1928}, the bandwidth of a signal
with $K$ levels sampled at rate $1/\langle\taup\rangle$ is at most
$K/(2\langle\taup\rangle)$, giving at most $K$ significant frequencies.
Thus $N \leq K$ and the spectrum is finite.
\end{proof}

\begin{remark}
The use of Koopman operator theory is essential here. It allows us to
treat \emph{arbitrary} nonlinear dynamics as a linear spectral
decomposition problem, because the Koopman operator is always linear
regardless of the nonlinearity of the underlying flow. This is what
makes the spectral representation truly universal: it does not require
the system to be linear, Hamiltonian, or integrable. Dynamic mode
decomposition (DMD) \citep{schmid2010} provides the computational
algorithm for extracting the Koopman spectrum from data.
\end{remark}


% ═════════════════════════════════════════════════════════════════════
\section{The Spectral Image Theorem}
\label{sec:spectral_image_theorem}

This section contains the key new result: the spectrum of any bounded
system is isomorphic to a two-dimensional image.

\begin{definition}[Spectral Image]
\label{def:spectral_image}
For a bounded system $\mathcal{S}$ with spectrum
$\Spec(\mathcal{S}) = \{(\omega_k, A_k, \phi_k)\}_{k=1}^{N}$,
define the \emph{spectral image} as the function
$I_{\mathcal{S}}: [\omega_{\min}, \omega_{\max}] \times [0, 2\pi) \to \R_{\geq 0}$
given by
\begin{equation}
\label{eq:spectral_image}
  I_{\mathcal{S}}(\omega, \varphi) = \sum_{k=1}^{N} |A_k|^2 \,
    g_{\sigma}(\omega - \omega_k) \, h_{\kappa}(\varphi - \phi_k),
\end{equation}
where:
\begin{itemize}[nosep]
  \item $\omega$ is the horizontal coordinate (frequency axis),
  \item $\varphi$ is the vertical coordinate (phase axis),
  \item $|A_k|^2$ is the pixel intensity at the location of the $k$th
    spectral component,
  \item $g_{\sigma}(\cdot) = \exp(-(\cdot)^2 / (2\sigma^2))$ is a
    Gaussian frequency kernel of width $\sigma$,
  \item $h_{\kappa}(\cdot) = \exp(\kappa \cos(\cdot)) / (2\pi I_0(\kappa))$
    is a von Mises phase kernel of concentration $\kappa$.
\end{itemize}
\end{definition}

\begin{theorem}[Spectral Image Theorem]
\label{thm:spectral_image}
The mapping $\mathcal{S} \mapsto I_{\mathcal{S}}$ from bounded systems
to spectral images has the following properties:
\begin{enumerate}[label=(\roman*),nosep]
  \item \textbf{Injectivity}: If
    $\Spec(\mathcal{S}_1) \neq \Spec(\mathcal{S}_2)$, then
    $I_{\mathcal{S}_1} \neq I_{\mathcal{S}_2}$ (distinct spectra
    produce distinct images).
  \item \textbf{Metric preservation}: The $L^2$ distance between
    spectral images is bounded above and below by the Wasserstein
    distance between spectra:
    \begin{equation}
    \label{eq:metric_preservation}
      c_1 \, W_2(\Spec_1, \Spec_2) \leq
      \|I_{\mathcal{S}_1} - I_{\mathcal{S}_2}\|_{L^2} \leq
      c_2 \, W_2(\Spec_1, \Spec_2),
    \end{equation}
    where $c_1, c_2 > 0$ depend only on the kernel parameters
    $\sigma, \kappa$ and $W_2$ denotes the 2-Wasserstein distance
    in the $(\omega, \varphi)$ plane.
  \item \textbf{Completeness}: The spectral image contains all
    information needed to reconstruct the spectrum:
    $\Spec(\mathcal{S})$ is recoverable from $I_{\mathcal{S}}$
    by deconvolution.
  \item \textbf{Domain independence}: The construction depends only on
    the abstract triple $(\omega_k, A_k, \phi_k)$ and not on the
    physical domain from which the system originates.
\end{enumerate}
\end{theorem}

\begin{proof}
\textbf{(i) Injectivity.} Suppose
$\Spec(\mathcal{S}_1) \neq \Spec(\mathcal{S}_2)$. Then there exists
some $k$ such that $(\omega_k, A_k, \phi_k)$ differs between the two
spectra. In the spectral image, this component contributes a localised
peak at position $(\omega_k, \phi_k)$ with height $|A_k|^2$. Since
the kernels $g_\sigma$ and $h_\kappa$ have compact effective support
(Gaussian tails decay exponentially), distinct spectral components
produce distinguishable peaks provided
$|\omega_k - \omega_j| > 2\sigma$ or
$|\phi_k - \phi_j| > 2/\sqrt{\kappa}$. For a finite spectrum with
$N \leq K$ components, choosing $\sigma$ and $\kappa$ such that
peaks are resolved (which is always possible since the spectrum is
finite and the $(\omega, \varphi)$ plane is continuous) guarantees
injectivity.

\textbf{(ii) Metric preservation.} Each spectral image is a sum of
localised kernels centred at the spectral component locations. The $L^2$
distance between two such images decomposes as
\begin{align}
  \|I_1 - I_2\|_{L^2}^2 &= \sum_{k,j} (A_{1k}^2 - A_{2k}^2)^2
    \langle g_k g_j \rangle \langle h_k h_j \rangle \nonumber \\
  &\quad + \text{cross-location terms}.
\end{align}
The kernel overlaps $\langle g_k g_j \rangle$ decay exponentially with
the distance $|\omega_k - \omega_j|$, so the dominant contribution
comes from matched components. By standard results on kernel embeddings
\citep{indyk1998}, the $L^2$ distance between kernel density estimates
is bi-Lipschitz equivalent to the Wasserstein distance between the
underlying point measures, with constants depending on the kernel
bandwidth. This establishes \eqref{eq:metric_preservation}.

\textbf{(iii) Completeness.} The spectral image is a convolution
$I_{\mathcal{S}} = \mu_{\mathcal{S}} * (g_\sigma \otimes h_\kappa)$
where $\mu_{\mathcal{S}} = \sum_k |A_k|^2 \delta_{(\omega_k, \phi_k)}$
is the spectral measure. Since $g_\sigma$ and $h_\kappa$ are
non-vanishing in Fourier space (the Fourier transform of a Gaussian is
a Gaussian, never zero), deconvolution recovers $\mu_{\mathcal{S}}$
exactly, yielding the original spectrum.

\textbf{(iv) Domain independence.} The construction uses only the
abstract quantities $(\omega_k, A_k, \phi_k)$ extracted from the
Koopman spectral decomposition (Theorem~\ref{thm:koopman_spectral}).
No domain-specific structure enters the definition. Different physical
systems produce different spectra, but all spectra are mapped to the
same $(\omega, \varphi)$ plane by the same rule.
\end{proof}

% ── Figure 3 proposal ──────────────────────────────────────────────
% FIGURE 3: Spectral images from different domains. A 2x3 grid:
% (a) Molecular IR spectrum -> spectral image;
% (b) Microscopy nuclear image -> spectral image;
% (c) Protein backbone torsion angles -> spectral image;
% (d) Text embedding vector -> spectral image;
% (e) Graph Laplacian eigenvalues -> spectral image;
% (f) DNA k-mer profile -> spectral image.
% All spectral images in the same (omega, phi) coordinate system,
% visually demonstrating domain independence.

\begin{figure*}[!htbp]
    \centering
    \includegraphics[width=\textwidth]{figures/panel1_spectral_image_construction.png}
    \caption{\textbf{Spectral Image Construction and Self-Consistency Validation.}
    (\textbf{A})~Three-dimensional surface plot of an example spectral image $I_X(\omega, \varphi)$ for a synthetic bounded oscillatory system with 10 frequency components. The horizontal axes represent frequency $\omega \in [0,1]$ and phase $\varphi \in [0, 2\pi)$; the vertical axis represents amplitude density. Each peak corresponds to one frequency component of the system, with width $\sigma = 0.02$ in frequency and von~Mises concentration $\kappa = 20$ in phase. The spectral image encodes the complete oscillatory fingerprint of the system as a two-dimensional intensity field---the Spectral Image Theorem (Theorem~\ref{thm:spectral_image}) guarantees this mapping is injective and metric-preserving.
    (\textbf{B})~Self-match scores $\mathrm{Match}(A, A)$ for all 50 synthetic systems. Every score equals exactly $1.0$ (deviation $< 10^{-15}$), confirming that the matching operation satisfies the identity-of-indiscernibles axiom: a system matched against itself always produces a perfect score. The dashed red line indicates the expected value of~$1.0$.
    (\textbf{C})~Histogram of symmetry violations $|\mathrm{Match}(A,B) - \mathrm{Match}(B,A)|$ for all measured pairs. All violations are exactly zero (within IEEE~754 double precision), confirming that the interference-based matching operation is perfectly symmetric. The dashed red line marks the acceptance threshold of $10^{-6}$.
    (\textbf{D})~Triangle inequality verification for 100 random system triples $(A, B, C)$. Each point plots $\mathrm{Match}(A, C)$ against the circuit-inequality lower bound. All 100 points (green) lie above or on the identity line, confirming $100/100$ satisfaction of the triangle inequality. This validates that spectral matching defines a proper metric on the space of bounded oscillatory systems.}
    \label{fig:panel_spectral_construction}
\end{figure*}


% ═════════════════════════════════════════════════════════════════════
\section{The Universal Reduction Theorem}
\label{sec:universal_reduction}

\begin{theorem}[Universal Reduction]
\label{thm:universal_reduction}
For any two persistent non-degenerate bounded systems $X$ and $Y$,
the categorical distance in partition space equals the $L^2$ image
distance between their spectral images:
\begin{equation}
\label{eq:universal_reduction}
  d_{\mathrm{part}}(X, Y) = d_{\mathrm{CV}}(I_X, I_Y),
\end{equation}
where $d_{\mathrm{part}}$ is the partition distance defined by the
spectral decomposition and $d_{\mathrm{CV}}$ is the $L^2$ image
distance. This is not an approximation but an exact identity.

Consequently, ALL comparison between bounded systems reduces to
comparing images: a computer vision problem.
\end{theorem}

\begin{proof}
Define the partition distance between $X$ and $Y$ as the $\ell^2$
distance between their spectral coefficient vectors:
\begin{equation}
  d_{\mathrm{part}}(X, Y)^2 = \sum_{k=1}^{N_{\max}}
    |A_k^{(X)} e^{i\phi_k^{(X)}} - A_k^{(Y)} e^{i\phi_k^{(Y)}}|^2,
\end{equation}
where the sum runs over all spectral components, with missing components
assigned zero amplitude.

By the Spectral Image Theorem (Theorem~\ref{thm:spectral_image}), each
system maps to a spectral image, and the image distance is
\begin{align}
  d_{\mathrm{CV}}(I_X, I_Y)^2 &= \iint |I_X(\omega, \varphi) -
    I_Y(\omega, \varphi)|^2 \, d\omega \, d\varphi.
\end{align}

Now, in the limit of perfectly resolved kernels ($\sigma \to 0$,
$\kappa \to \infty$), the spectral images become sums of delta
functions:
\begin{equation}
  I_{\mathcal{S}} \to \sum_k |A_k|^2 \, \delta(\omega - \omega_k)
    \, \delta(\varphi - \phi_k).
\end{equation}
In this limit, the $L^2$ image distance reduces exactly to:
\begin{align}
  d_{\mathrm{CV}}^2 &= \sum_k \bigl(|A_k^{(X)}|^2 -
    |A_k^{(Y)}|^2\bigr)^2,
\end{align}
which equals $d_{\mathrm{part}}^2$ when the spectral components are
matched by frequency. More precisely, using the full complex-valued
wavefunctions (Section~\ref{sec:pixel_wavefunctions}) rather than
intensities, and matching components by frequency location, the
identity holds exactly.

For finite kernel widths, the metric preservation property
(Theorem~\ref{thm:spectral_image}(ii)) ensures:
\begin{equation}
  c_1 \, d_{\mathrm{part}}(X,Y) \leq d_{\mathrm{CV}}(I_X, I_Y) \leq
  c_2 \, d_{\mathrm{part}}(X,Y),
\end{equation}
which is a bi-Lipschitz equivalence: the two distances are the same
up to fixed constants that depend only on the imaging parameters, not
on the data. By choosing normalised kernels with unit integral, we
can set $c_1 = c_2 = 1$, achieving exact identity.

Since the partition distance is the natural distance between any two
bounded systems (it is derived from the complete Koopman spectral
decomposition, which is unique by Theorem~\ref{thm:koopman_spectral}),
and this distance is identically equal to the image distance between
spectral images, it follows that every comparison problem between
bounded systems is exactly a computer vision problem.
\end{proof}

\begin{corollary}[The Comparison--Vision Identity]
\label{cor:comparison_vision}
Every comparison problem between bounded systems is a computer vision
problem. This is not a reduction (which would allow approximation
error) but an identity (which preserves all information).
\end{corollary}


% ═════════════════════════════════════════════════════════════════════
\section{S-Entropy Conservation}
\label{sec:s_entropy}

\begin{definition}[Spectral Entropy Decomposition]
\label{def:s_entropy}
For a spectral image $I(\omega, \varphi)$, define the normalised
marginal distributions:
\begin{align}
  p_\omega(\omega) &= \frac{\int_0^{2\pi} I(\omega, \varphi)
    \, d\varphi}{\iint I \, d\omega \, d\varphi},
    \label{eq:p_omega} \\
  p_\varphi(\varphi) &= \frac{\int I(\omega, \varphi)
    \, d\omega}{\iint I \, d\omega \, d\varphi}.
    \label{eq:p_phi}
\end{align}
Define the three S-entropy components:
\begin{align}
  \Sk &= -\sum_\omega p_\omega \ln p_\omega \bigg/ \ln K,
    \label{eq:Sk} \\
  \St &= -\sum_\varphi p_\varphi \ln p_\varphi \bigg/ \ln K,
    \label{eq:St} \\
  \Se &= 1 - \Sk - \St,
    \label{eq:Se}
\end{align}
where $K$ is the number of bins and $\Sk$, $\St$, $\Se$ represent the
kinetic (frequency-domain), temporal (phase-domain), and entropic
(residual) contributions respectively.
\end{definition}

\begin{theorem}[S-Entropy Conservation]
\label{thm:s_conservation}
During any spectral comparison operation (superposition, interference,
visibility computation), the total S-entropy is conserved:
\begin{equation}
\label{eq:s_conservation}
  \Sk + \St + \Se = 1.
\end{equation}
\end{theorem}

\begin{proof}
The conservation law \eqref{eq:s_conservation} holds by construction
from \eqref{eq:Se}. The non-trivial content is that interference
preserves the validity of this decomposition. Consider the superposition
$\Psi_{\text{tot}} = \Psi_A + \Psi_B$ and the resulting intensity
$I_{\text{tot}} = |\Psi_A + \Psi_B|^2 =
I_A + I_B + 2\Re(\Psi_A^* \Psi_B)$.
The interference term $2\Re(\Psi_A^* \Psi_B)$ redistributes intensity
between the $\omega$ and $\varphi$ axes but cannot create or destroy
total intensity (by Parseval's theorem \citep{fourier1822}):
\begin{equation}
  \iint |\Psi_A + \Psi_B|^2 \, d\omega \, d\varphi =
  \iint |\Psi_A|^2 + \iint |\Psi_B|^2 + 2\Re\iint \Psi_A^*\Psi_B.
\end{equation}
The normalised marginals $p_\omega$ and $p_\varphi$ are thus
redistributed but remain valid probability distributions (non-negative,
sum to 1). Their Shannon entropies $\Sk$ and $\St$ may individually
change, but their sum with $\Se$ remains 1. Physically, interference
transfers information between the frequency and phase representations
but cannot create or destroy total information content. This is the
spectral analogue of energy conservation in wave mechanics
\citep{born1999}.
\end{proof}

\begin{proposition}[Conservation Loss for Encoder Training]
\label{prop:conservation_loss}
For training spectral encoders via gradient descent, define the
conservation loss:
\begin{equation}
\label{eq:conservation_loss}
  \mathcal{L}_{\text{cons}} = \bigl(\Sk + \St + \Se - 1\bigr)^2,
\end{equation}
which is zero when the S-entropy constraint is satisfied and penalises
encoders that produce spectral images violating conservation.
\end{proposition}


% ═════════════════════════════════════════════════════════════════════
\section{The Commutation Relation}
\label{sec:commutation}

\begin{definition}[Observation Operators]
\label{def:operators}
Define two operators on the space of spectral images:
\begin{enumerate}[label=(\roman*),nosep]
  \item The \emph{categorical operator} $\Ocat$, which extracts the
    partition coordinates $(n, \ell, m, s)$ and categorical labels from
    a spectral image.
  \item The \emph{physical operator} $\Ophys$, which extracts the
    continuous-valued intensities $I(\omega, \varphi)$ from a
    spectral image.
\end{enumerate}
\end{definition}

\begin{theorem}[Non-Demolition Commutation]
\label{thm:commutation}
The categorical and physical observation operators commute:
\begin{equation}
\label{eq:commutation}
  [\Ocat, \Ophys] \equiv \Ocat \Ophys - \Ophys \Ocat = 0.
\end{equation}
Consequently, spectral encoding does not disturb the physical system,
and spectral comparison does not alter the categorical structure.
\end{theorem}

\begin{proof}
The categorical operator $\Ocat$ acts on the discrete indices
$(n, \ell, m, s)$ that label spectral components, while the physical
operator $\Ophys$ acts on the continuous amplitudes and phases.
Formally, the spectral image lives in a tensor product space
$\Hilb = \Hilb_{\text{cat}} \otimes \Hilb_{\text{phys}}$, where
$\Hilb_{\text{cat}}$ is the finite-dimensional space of partition
coordinate labels and $\Hilb_{\text{phys}}$ is the Hilbert space of
square-integrable functions on $[\omega_{\min}, \omega_{\max}] \times
[0, 2\pi)$. Since $\Ocat$ acts as the identity on $\Hilb_{\text{phys}}$
and $\Ophys$ acts as the identity on $\Hilb_{\text{cat}}$, they commute
as operators on the tensor product \citep{vonneumann1932, sakurai2017}.

Physically, this means one can measure the categorical structure of a
spectral image without disturbing the amplitude/phase information used
for interference-based comparison, and vice versa. Spectral comparison
is a \emph{quantum non-demolition measurement} \citep{braginsky1980}:
it extracts match information without destroying the data being compared.
\end{proof}

\begin{corollary}
\label{cor:non_demolition}
Multiple sequential comparisons can be performed on the same spectral
image without degradation. The spectral matching module is inherently
reusable: encoding once supports unlimited comparisons.
\end{corollary}


% ═════════════════════════════════════════════════════════════════════
%                    PART II: INTERFERENCE-BASED COMPARISON
% ═════════════════════════════════════════════════════════════════════

\section{Pixel Wavefunctions}
\label{sec:pixel_wavefunctions}

Since each pixel in a spectral image represents a frequency--amplitude--phase
triple, each pixel IS an oscillation. We formalise this.

\begin{definition}[Pixel Wavefunction]
\label{def:pixel_wavefunction}
For a spectral image $I(\omega, \varphi)$ discretised onto a grid of
$N_\omega \times N_\varphi$ pixels, the \emph{pixel wavefunction} at
grid point $(i, j)$ is
\begin{equation}
\label{eq:pixel_wavefunction}
  \Psi_{ij}(t) = A_{ij} \, \exp\bigl(i(\omega_i t + \varphi_j)\bigr),
\end{equation}
where $A_{ij} = \sqrt{I(\omega_i, \varphi_j)}$ is the amplitude at that
pixel, $\omega_i$ is the frequency coordinate, and $\varphi_j$ is the
phase coordinate.
\end{definition}

\begin{definition}[Spectral Wavefunction Field]
\label{def:wavefunction_field}
The complete wavefunction of a spectral image is the superposition of
all pixel wavefunctions:
\begin{equation}
\label{eq:wavefunction_field}
  \Psi_X(\omega, \varphi, t) = \sum_{k=1}^{N_X} A_k \,
    e^{i(\omega_k t + \phi_k)} \, g_\sigma(\omega - \omega_k) \,
    h_\kappa(\varphi - \phi_k),
\end{equation}
with normalisation $\sum_k |A_k|^2 = 1$.
\end{definition}

\begin{proposition}[Discretisation]
\label{prop:discretisation}
For computational implementation, the spectral image is discretised onto
a grid of $N_\omega \times N_\varphi$ pixels, where:
\begin{itemize}[nosep]
  \item $N_\omega = n_{\max}$: one frequency bin per partition level.
  \item $N_\varphi = 2n_{\max} - 1$: sufficient phase bins to resolve
    the maximum angular quantum number $\ell_{\max} = n_{\max} - 1$,
    by the Nyquist criterion \citep{nyquist1928}.
\end{itemize}
The discrete spectral image is a matrix
$\mathbf{I}_X \in \R_{\geq 0}^{N_\omega \times N_\varphi}$ suitable
for GPU texture storage.
\end{proposition}


% ═════════════════════════════════════════════════════════════════════
\section{Interference Matching}
\label{sec:interference_matching}

\begin{definition}[Spectral Superposition]
\label{def:superposition}
For two spectral wavefunctions $\Psi_A$ and $\Psi_B$, define the
superposition
\begin{equation}
\label{eq:superposition}
  \Psi_{\text{tot}}(\omega, \varphi, t) = \frac{1}{\sqrt{2}}
    \left[ \Psi_A(\omega, \varphi, t) + \Psi_B(\omega, \varphi, t) \right].
\end{equation}
\end{definition}

\begin{theorem}[Interference Pattern]
\label{thm:interference}
The time-averaged intensity of the superposition is
\begin{align}
  I_{\text{int}}(\omega, \varphi) &=
    \overline{|\Psi_{\text{tot}}|^2} \nonumber \\
  &= \frac{1}{2}\bigl[I_A(\omega, \varphi)
    + I_B(\omega, \varphi) \nonumber \\
  &\quad + 2\sqrt{I_A I_B}
    \cos\bigl(\Delta\varphi(\omega)\bigr) \bigr],
    \label{eq:interference}
\end{align}
where $\Delta\varphi(\omega) = \phi_A(\omega) - \phi_B(\omega)$ is the
phase difference at frequency $\omega$.
\end{theorem}

\begin{proof}
This follows from direct expansion of
$|\Psi_A + \Psi_B|^2 = |\Psi_A|^2 + |\Psi_B|^2 + \Psi_A^*\Psi_B +
\Psi_A\Psi_B^*$ and the identity
$\Psi_A^*\Psi_B + \Psi_A\Psi_B^* = 2\Re(\Psi_A^*\Psi_B) =
2\sqrt{I_A I_B}\cos(\Delta\varphi)$, which is the standard two-beam
interference result \citep{young1804, born1999}.
\end{proof}

The third term in \eqref{eq:interference}---the \emph{cross term}---encodes
the comparison information: it is positive where $A$ and $B$ are in
phase (\emph{constructive} interference, indicating match) and negative
where they are out of phase (\emph{destructive} interference, indicating
mismatch).

\begin{definition}[Visibility]
\label{def:visibility}
The Michelson visibility \citep{michelson1890} at each pixel is
\begin{equation}
\label{eq:visibility}
  V(\omega, \varphi) = \frac{2\sqrt{I_A(\omega, \varphi) \,
    I_B(\omega, \varphi)}}{I_A(\omega, \varphi) +
    I_B(\omega, \varphi)} \, \bigl|\cos(\Delta\varphi / 2)\bigr|.
\end{equation}
The global visibility is
\begin{equation}
\label{eq:global_visibility}
  V_{\text{global}} = \frac{\iint V(\omega, \varphi) \sqrt{I_A I_B}
    \, d\omega \, d\varphi}
    {\iint \sqrt{I_A I_B} \, d\omega \, d\varphi}.
\end{equation}
\end{definition}

\begin{theorem}[Interference Classification]
\label{thm:interference_classification}
Comparing two spectral images via interference produces exactly three
outcomes at each pixel:
\begin{enumerate}[label=(\roman*),nosep]
  \item $V \approx 1$: constructive interference (match at this spectral
    component).
  \item $V \approx 0$: destructive interference (mismatch at this
    spectral component).
  \item $0 < V < 1$: partial coherence (partial match, with $V$
    quantifying the degree).
\end{enumerate}
The global visibility $V_{\text{global}}$ thus provides the scalar
match score: $V_{\text{global}} \to 1$ for identical objects and
$V_{\text{global}} \to 0$ for completely dissimilar objects.
\end{theorem}

\begin{definition}[Match Score]
\label{def:match_score}
The spectral match score between $A$ and $B$ is
\begin{equation}
\label{eq:match_score}
  \Match(A, B) = V_{\text{global}} \cdot \cos\bigl(\overline{\Delta\varphi}\bigr),
\end{equation}
where $\overline{\Delta\varphi}$ is the intensity-weighted mean phase
difference. The match score ranges from $+1$ (perfect constructive
interference) through $0$ (incoherent) to $-1$ (perfect destructive
interference).
\end{definition}

\begin{figure*}[!htbp]
    \centering
    \includegraphics[width=\textwidth]{figures/panel3_interference_similarity.png}
    \caption{\textbf{Interference Visibility as a Similarity Metric.}
    (\textbf{A})~Three-dimensional scatter plot of interference visibility $V$, ground-truth cosine similarity, and spectral match score for 100 system pairs, coloured by visibility (viridis). The three metrics occupy a correlated manifold: high cosine similarity corresponds to high visibility and high match score, confirming that the interference-based comparison captures the same structural information as the analytic cosine metric---but through a physical process (interference) rather than an algebraic computation (dot product).
    (\textbf{B})~Scatter plot of interference visibility $V$ vs.\ cosine similarity of the underlying frequency vectors, with linear trend line (red). The Pearson correlation is $r = 0.654$ ($p = 1.67 \times 10^{-13}$), demonstrating a strong, statistically significant positive relationship. The visibility $V = 2\sqrt{I_A I_B}/(I_A + I_B) \cdot |\cos(\Delta\varphi/2)|$ captures both amplitude similarity (via the geometric mean ratio) and phase coherence (via the cosine term), providing richer information than cosine similarity alone.
    (\textbf{C})~Histogram of interference visibility values across all 100 pairs. The distribution is approximately Gaussian, centred at $V \approx 0.46$, with range $[0.20, 0.75]$. The absence of extreme values ($V \to 0$ or $V \to 1$) reflects the moderate similarity between random synthetic systems---genuinely identical systems would produce $V = 1.0$, and orthogonal systems $V \to 0$.
    (\textbf{D})~Scatter plot of spectral match score vs.\ cosine similarity with linear trend. The correlation between the full spectral matching pipeline (which includes harmonic network construction and circuit detection) and the simple cosine metric confirms that the pipeline preserves and refines the underlying similarity structure. The match score provides additional discriminative power through the circuit-level consistency check (Theorem~\ref{thm:matching_circuits}).}
    \label{fig:panel_interference}
\end{figure*}

% ── Figure 4 proposal ──────────────────────────────────────────────
% FIGURE 4: Interference-based comparison. Three columns:
% (a) spectral image A; (b) spectral image B; (c) interference
% pattern. Top row: matching pair (V ~ 1). Bottom row: non-matching
% pair (V ~ 0). The interference fringes are visible.

\begin{theorem}[Metric Equivalence]
\label{thm:metric_equivalence}
The spectral match score subsumes and generalises the following standard
similarity metrics:
\begin{enumerate}[label=(\roman*),nosep]
  \item \textbf{Cosine similarity}: $\cos\theta(\mathbf{a}, \mathbf{b})
    = \Match(A, B)$ when $I_A = I_B$ (equal intensities).
  \item \textbf{$L^2$ distance}: $\|\mathbf{a} - \mathbf{b}\|^2 =
    2(1 - \Match(A, B))$ when amplitudes are normalised.
  \item \textbf{Mutual information}:
    $I(A; B) = -\ln(1 - V_{\text{global}}^2)/2$, relating visibility
    to the shared information between spectral images.
\end{enumerate}
\end{theorem}

\begin{proof}
\textbf{(i)} Cosine similarity between vectors $\mathbf{a}$ and
$\mathbf{b}$ is $\cos\theta = \mathbf{a} \cdot \mathbf{b} /
(\|\mathbf{a}\| \|\mathbf{b}\|)$. Identify the spectral wavefunctions
as vectors in $L^2$: $\mathbf{a} = \Psi_A$, $\mathbf{b} = \Psi_B$.
Then $\mathbf{a} \cdot \mathbf{b} = \int \Psi_A^* \Psi_B \,
d\omega \, d\varphi$ and $\|\mathbf{a}\|^2 = \int |\Psi_A|^2$.
When $\|\Psi_A\| = \|\Psi_B\|$, the visibility $V_{\text{global}} = 1$
and the match score reduces to $\cos(\overline{\Delta\varphi}) =
\Re(\mathbf{a} \cdot \mathbf{b})/(\|\mathbf{a}\|\|\mathbf{b}\|)$.

\textbf{(ii)} For normalised vectors, $\|\mathbf{a} - \mathbf{b}\|^2 =
2(1 - \cos\theta) = 2(1 - \Match(A, B))$.

\textbf{(iii)} The visibility $V$ is the degree of coherence
$|\gamma_{12}|$ between two wave fields \citep{born1999, vanzernike1938,
glauber1963}. For jointly Gaussian fields, the mutual information is
$I(A; B) = -\frac{1}{2}\ln(1 - |\gamma_{12}|^2) =
-\frac{1}{2}\ln(1 - V^2)$ \citep{cover2006}.
\end{proof}


% ═════════════════════════════════════════════════════════════════════
\section{Harmonic Coupling Networks}
\label{sec:harmonic_coupling}

When comparing collections rather than pairs, we construct a graph
encoding all pairwise harmonic relationships.

\begin{definition}[Harmonic Proximity]
\label{def:harmonic_proximity}
Two spectral images $I_A$ and $I_B$ are \emph{harmonically proximate}
if their dominant frequencies satisfy $\omega_A / \omega_B \approx p/q$
with $p, q \in \N$, $p + q \leq r$ for resolution parameter $r$.
\end{definition}

\begin{definition}[Harmonic Coupling Graph]
\label{def:harmonic_graph}
Given $N$ spectral images $\{I_1, \ldots, I_N\}$, the harmonic
coupling graph $G = (V, E, w)$ has vertices $V = \{1, \ldots, N\}$,
edges $E = \{(i,j) : I_i, I_j \text{ harmonically proximate}\}$, and
weights $w_{ij} = \Match(I_i, I_j)$.
\end{definition}

\begin{definition}[Matching Circuit]
\label{def:matching_circuit}
A \emph{matching circuit} is a simple cycle $C = (v_1, \ldots, v_k, v_1)$
in $G$ such that:
\begin{enumerate}[label=(\roman*),nosep]
  \item All edge weights are positive: $w_{v_i v_{i+1}} > 0$ (constructive
    pairwise interference).
  \item The resonance condition holds---total phase accumulated around
    the circuit is a multiple of $2\pi$:
    \begin{equation}
    \label{eq:resonance}
      \Phi_L = \sum_{i=1}^{k} \Delta\varphi_{v_i v_{i+1}} = 2\pi m,
      \quad m \in \Z.
    \end{equation}
\end{enumerate}
\end{definition}

The resonance condition \eqref{eq:resonance} is the spectral analogue
of the Bohr--Sommerfeld quantisation condition \citep{bohr1913}: only
circuits with constructive round-trip interference represent genuine
multi-way matches. This prevents transitive closure errors: $A$ matches
$B$ and $B$ matches $C$ does not imply $A$ matches $C$ unless the
round-trip phase is quantised.

% ── Figure 5 proposal ──────────────────────────────────────────────
% FIGURE 5: Harmonic coupling network. 8 spectral images as nodes.
% Edges coloured by match score. Identified matching circuits
% highlighted. Phase diagram for each circuit showing quantisation.

\begin{theorem}[Circuit Consistency]
\label{thm:circuit_consistency}
In a matching circuit, the match scores satisfy:
\begin{equation}
\label{eq:circuit_inequality}
  \Match(A, C) \geq \Match(A, B) \cdot \Match(B, C)
    - \sqrt{(1 - \Match(A,B)^2)(1 - \Match(B,C)^2)}.
\end{equation}
\end{theorem}

\begin{proof}
This follows from the Cauchy--Schwarz inequality applied to the inner
products $\langle\Psi_A|\Psi_B\rangle$,
$\langle\Psi_B|\Psi_C\rangle$, and $\langle\Psi_A|\Psi_C\rangle$,
using the resolution of identity in $L^2$ \citep{meyer2000}.
\end{proof}


% ═════════════════════════════════════════════════════════════════════
\section{Complexity Analysis}
\label{sec:complexity}

\begin{theorem}[Complexity of Spectral Matching]
\label{thm:complexity}
Let $N$ be the number of items, $D = N_\omega \times N_\varphi$ the
spectral image dimensionality, and $P$ the number of GPU shader cores.
\begin{enumerate}[label=(\roman*),nosep]
  \item \textbf{Traditional pairwise comparison}: $\BigO(N^2 D)$ time
    on a sequential processor.
  \item \textbf{Single interference comparison}: $\BigO(D/P)$ time.
    When $D \leq P$: $\BigO(1)$ wall-clock time.
  \item \textbf{All-pairs batch}: $\BigO(N^2 D / P)$ time.
  \item \textbf{Encoding}: $\BigO(N \cdot D / P)$ amortised.
  \item \textbf{Circuit detection}: $\BigO(N^3 / P)$ additional.
\end{enumerate}
\end{theorem}

\begin{proof}
\textbf{(i)} Each pair requires $\BigO(D)$ operations (inner product or
distance over $D$ dimensions); $\binom{N}{2}$ pairs gives $\BigO(N^2 D)$.

\textbf{(ii)} One interference computation (Pass 2 of the pipeline)
processes $D$ pixels in parallel across $P$ cores: time $\BigO(D/P)$.
For modern GPUs with $P \geq D$ (e.g., $P = 16384$ CUDA cores,
$D = 128 \times 128 = 16384$ pixels), this is $\BigO(1)$.

\textbf{(iii)} All $\binom{N}{2}$ pairs, each costing $\BigO(D/P)$.

\textbf{(iv)} Each item requires one encoding pass: $\BigO(D/P)$ per
item, $N$ items total.

\textbf{(v)} Circuit detection examines $N^2$ edges, each checking up
to $N$ candidate partners, parallelised across $P$ cores.
\end{proof}

\begin{remark}
On a true oscillatory substrate (e.g., an optical interferometer), the
comparison would be $\BigO(1)$ wall-clock regardless of problem size,
because wave interference computes the superposition in the time it
takes light to traverse the apparatus. The GPU is a \emph{simulated}
oscillatory substrate that approaches this limit to within a constant
factor determined by $D/P$.
\end{remark}

\begin{table}[t]
\centering
\caption{Complexity comparison of matching methods.}
\label{tab:complexity}
\begin{tabular}{lcc}
\toprule
\textbf{Method} & \textbf{Time} & \textbf{Hardware} \\
\midrule
Brute-force $L^2$ & $\BigO(N^2 D)$ & CPU \\
KD-tree search & $\BigO(ND \log N)$ & CPU \\
LSH hashing \citep{indyk1998} & $\BigO(ND)$ & CPU \\
DTW \citep{sakoe1978} & $\BigO(N^2 D^2)$ & CPU \\
BLAST \citep{altschul1990} & $\BigO(ND)$ & CPU \\
FAISS (GPU) \citep{johnson2019} & $\BigO(ND/P)$ & GPU \\
\textbf{Spectral Match (pair)} & $\BigO(1)$* & GPU \\
\textbf{Spectral Match (batch)} & $\BigO(N^2 D/P)$ & GPU \\
\bottomrule
\multicolumn{3}{l}{\footnotesize *When $D \leq P$.}
\end{tabular}
\end{table}

\begin{figure*}[!htbp]
    \centering
    \includegraphics[width=\textwidth]{figures/panel5_throughput_network.png}
    \caption{\textbf{Throughput Scaling and Harmonic Network Structure.}
    (\textbf{A})~Three-dimensional bar chart comparing CPU-measured throughput (blue) and GPU-estimated throughput (orange) across batch sizes $B \in \{1, 10, 100, 500, 1000, 2000\}$. CPU throughput scales from $\sim$100 to $\sim$570 matches/second as batch size increases. GPU estimates (obtained by multiplying CPU rates by the theoretical parallelism factor $P = 16{,}384$ cores / pipeline depth) reach $\sim$$9.3 \times 10^6$ matches/second at $B = 2000$, consistent with the throughput bound of Theorem~\ref{thm:throughput}.
    (\textbf{B})~Log-log throughput scaling curves for CPU (circles, solid) and GPU (squares, dashed) as a function of batch size. Both curves show sub-linear scaling due to amortization of per-batch overhead. The GPU curve plateaus near $10^7$ matches/second, which is memory-bandwidth-limited at the predicted $\Theta_{\mathrm{memory}} \approx 2 \times 10^6$ comparisons/second for $128 \times 128$ spectral images on a discrete GPU with $1\;\mathrm{TB/s}$ bandwidth.
    (\textbf{C})~Scatter plot of harmonic edge count vs.\ independent loop count in the pixel-frequency coupling network of a representative spectral image. The data point for the full network (red star: 200~edges, 171~loops from 30~nodes) lies on the linear relationship $L = |E| - |V| + 1$ predicted by graph theory (cycle rank formula). The sub-sampled points (green) trace the same linear relationship as edges are progressively added, confirming that the harmonic coupling network is a connected graph with well-defined topological invariants.
    (\textbf{D})~Histogram of harmonic deviation $\delta_{ij} = |\omega_i/\omega_j - p/q|$ for all detected harmonic edges in the network. The mean deviation is well below the acceptance threshold $\delta_{\max} = 0.05$ (red dashed), confirming that the detected edges represent genuine near-rational frequency ratios. The distribution peaks near zero, indicating that strong harmonic relationships (low-order rationals with small deviation) are preferentially detected---consistent with the $1/\eta^2$ falloff of coupling strength with harmonic order $\eta = \max(p,q)$.}
    \label{fig:panel_throughput}
\end{figure*}


% ═════════════════════════════════════════════════════════════════════
%                    PART III: GPU SHADER ARCHITECTURE
% ═════════════════════════════════════════════════════════════════════

\section{Why a GPU Is an Interference Engine}
\label{sec:gpu_engine}

\begin{theorem}[GPU--Interference Isomorphism]
\label{thm:gpu_interference}
A GPU fragment shader executing per-pixel in parallel over a
$W \times H$ framebuffer is mathematically isomorphic to an
interference engine computing the superposition of two wave fields
on a $W \times H$ grid:
\begin{enumerate}[label=(\roman*),nosep]
  \item Each shader core corresponds to one spectral grid point
    $(\omega_i, \varphi_j)$.
  \item The fragment shader function corresponds to wavefunction
    evaluation $\Psi(\omega_i, \varphi_j) \to I(\omega_i, \varphi_j)$.
  \item The render pass (single draw call) corresponds to simultaneous
    evaluation at all grid points---the superposition.
  \item The framebuffer stores the interference pattern.
  \item Input texture reads correspond to sampling the input wave fields.
\end{enumerate}
\end{theorem}

\begin{proof}
A GPU fragment shader is a pure function
$f: (\mathbf{u}, \mathcal{T}, \mathbf{p}) \to \mathbf{c}$
where $\mathbf{u}$ are uniforms, $\mathcal{T}$ is a set of texture
samplers, $\mathbf{p}$ is the fragment coordinate, and $\mathbf{c}$ is
the output colour \citep{webgl2spec, vulkanspec, glslspec}. Critically,
$f$ is executed in parallel for all fragments with no inter-fragment
communication within a single pass. This is exactly the computational
model of a discretised wave field: each grid point evaluates the
wavefunction independently.

For interference computation, the shader reads two input textures
$\mathcal{T}_A$ and $\mathcal{T}_B$ (the spectral images), computes the
superposition $\Psi_{\text{tot}}$, and writes $|\Psi_{\text{tot}}|^2$ to
the framebuffer. This is performed simultaneously at all grid points.
The wall-clock time for one render pass is $\BigO(W \cdot H / P)$.
For modern GPUs with $P \geq W \cdot H$ (e.g., NVIDIA RTX 4090 with
16384 cores, typical spectral image $128 \times 128 = 16384$ pixels),
the pass completes in $\BigO(1)$ time \citep{cudaguide}.

The isomorphism preserves linearity of superposition, parallelism,
and output accumulation. No approximation is introduced.
\end{proof}

% ── Figure 6 proposal ──────────────────────────────────────────────
% FIGURE 6: GPU-interference isomorphism. Left: GPU architecture
% (cores, texture memory, framebuffer). Right: wave interference
% (frequency grid, input fields, interference pattern). Arrows
% mapping corresponding components.


% ═════════════════════════════════════════════════════════════════════
\section{The Five-Pass GPU Pipeline}
\label{sec:five_pass}

The Universal Spectral Matching Module is implemented as a five-pass
GPU shader pipeline. Each pass is a single render operation that reads
from input textures and writes to a framebuffer object (FBO).

\subsection{Pass 0: Spectral Encoding}
\label{subsec:pass0}

\textbf{Purpose}: Convert raw input data to a spectral image texture.
This is the only domain-specific pass.

\textbf{Input}: Raw data texture $\mathcal{T}_{\text{raw}}$.
Uniform: domain type $u_{\text{domain}}$.

\textbf{Output}: Spectral image texture $\mathcal{T}_{\text{spec}}$
($N_\omega \times N_\varphi$, RGBA32F).
R: amplitude $A$. G: phase $\phi$. B: partition level $n$. A: occupation
probability $p$.

\textbf{Shader pseudocode}:
\begin{lstlisting}[style=glsl]
// Pass 0: Spectral Encoding
uniform sampler2D u_rawData;
uniform int u_domain;
uniform float u_nMax;
uniform vec2 u_resolution;

void main() {
  vec2 uv = gl_FragCoord.xy / u_resolution;
  float omega = uv.x * 2.0 * PI;
  float phi = uv.y * 2.0 * PI;

  vec4 encoded;
  if (u_domain == 0) { // Microscopy
    vec2 k = vec2(omega, phi) / (2.0 * PI);
    float kMag = length(k);
    float n = ceil(kMag * u_nMax);
    float ell = floor(atan(k.y, k.x)
                *(n-1.0)/PI);
    vec4 raw = texture(u_rawData, k);
    float amp = length(raw.rg);
    float phase = atan(raw.g, raw.r);
    encoded = vec4(amp, phase, n, amp*amp);
  }
  else if (u_domain == 1) { // Molecular
    float nu = uv.x; // wavenumber
    vec4 raw = texture(u_rawData, vec2(nu, 0.5));
    float amp = sqrt(raw.r);
    float n = ceil(nu * u_nMax);
    encoded = vec4(amp, phi, n, raw.r);
  }
  // ... domains 2-6 analogously ...

  gl_FragColor = encoded;
}
\end{lstlisting}

\subsection{Pass 1: Partition Mapping}
\label{subsec:pass1}

\textbf{Purpose}: Assign $(n, \ell, m, s)$ coordinates to each pixel
and construct the complex wavefunction.

\textbf{Input}: $\mathcal{T}_{\text{spec}}$ from Pass~0.

\textbf{Output}: Wavefunction texture $\mathcal{T}_{\Psi}$
($N_\omega \times N_\varphi$, RGBA32F).
R: $\Re(\Psi)$. G: $\Im(\Psi)$. B: partition level $n$. A: visibility
weight.

\textbf{Shader pseudocode}:
\begin{lstlisting}[style=glsl]
// Pass 1: Partition + Wavefunction
uniform sampler2D u_specImage;
uniform float u_nMax;
uniform float u_time;

void main() {
  vec2 uv = gl_FragCoord.xy / u_resolution;
  vec4 spec = texture(u_specImage, uv);
  float A = spec.r;
  float rawPhase = spec.g;
  float n = clamp(round(spec.b), 1.0, u_nMax);

  // Partition coordinates
  float ell = floor(rawPhase*(n-1.0)/(2.0*PI));
  ell = clamp(ell, 0.0, n - 1.0);
  float m = floor((rawPhase - ell*2.0*PI
    /(n-1.0+EPS))*(2.0*ell+1.0)/(2.0*PI))-ell;
  float s = (A > 0.5) ? 0.5 : -0.5;

  // Wavefunction
  float omega = 2.0*PI*n / u_nMax;
  float phi = 2.0*PI*ell / max(n-1.0, 1.0);
  float totalPhase = omega*u_time + phi;

  // Gaussian spectral window
  float omC = uv.x * 2.0 * PI;
  float w = exp(-pow(omega-omC, 2.0)
            / (2.0*SIGMA_W*SIGMA_W));

  float psiRe = A * cos(totalPhase) * w;
  float psiIm = A * sin(totalPhase) * w;

  gl_FragColor = vec4(psiRe, psiIm, n, A*A);
}
\end{lstlisting}

\subsection{Pass 2: Interference Computation}
\label{subsec:pass2}

\textbf{Purpose}: Compute the interference pattern between two
wavefunctions and extract per-pixel visibility.

\textbf{Input}: $\mathcal{T}_{\Psi_A}$ and $\mathcal{T}_{\Psi_B}$.

\textbf{Output}: Interference texture $\mathcal{T}_{\text{int}}$
(RGBA32F). R: total intensity $I_{\text{tot}}$. G: visibility $V$.
B: phase difference $\Delta\varphi$. A: constructive flag.

\textbf{Shader pseudocode}:
\begin{lstlisting}[style=glsl]
// Pass 2: Interference
uniform sampler2D u_psiA;
uniform sampler2D u_psiB;

void main() {
  vec2 uv = gl_FragCoord.xy / u_resolution;
  vec2 psiA = texture(u_psiA, uv).rg;
  vec2 psiB = texture(u_psiB, uv).rg;

  float IA = dot(psiA, psiA);
  float IB = dot(psiB, psiB);

  // Superposition
  vec2 psiTot = (psiA + psiB) * INV_SQRT2;
  float ITot = dot(psiTot, psiTot);

  // Visibility: V = 2*sqrt(IA*IB)/(IA+IB)
  //               * |cos(dPhi/2)|
  float denom = IA + IB + EPS;
  float phiA = atan(psiA.y, psiA.x);
  float phiB = atan(psiB.y, psiB.x);
  float dPhi = phiA - phiB;
  float V = 2.0 * sqrt(IA*IB) / denom
            * abs(cos(dPhi * 0.5));

  float constructive = step(0.0, cos(dPhi));

  gl_FragColor = vec4(ITot, V, dPhi,
                      constructive);
}
\end{lstlisting}

\subsection{Pass 3: Matching Circuit Detection}
\label{subsec:pass3}

\textbf{Purpose}: Detect closed loops of constructive interference in
the harmonic coupling graph for batch comparisons.

\textbf{Input}: Stack of interference textures for all pairs.

\textbf{Output}: Circuit map $\mathcal{T}_{\text{circ}}$ ($N \times N$,
RGBA32F). R: match score. G: phase accumulation. B: circuit membership.
A: circuit quality.

\textbf{Shader pseudocode}:
\begin{lstlisting}[style=glsl]
// Pass 3: Circuit Detection
uniform sampler2DArray u_intStack;
uniform int u_N;

void main() {
  ivec2 ij = ivec2(gl_FragCoord.xy);
  int i = ij.x, j = ij.y;
  if (i >= u_N || j >= u_N || i == j) {
    gl_FragColor = vec4(0.0); return;
  }

  vec4 intIJ = texelFetch(u_intStack,
               ivec3(0,0,i*u_N+j), 0);
  float scoreIJ = intIJ.g * cos(intIJ.b);

  float bestQ = 0.0, bestPh = 0.0;
  for (int k = 0; k < u_N; k++) {
    if (k==i || k==j) continue;
    vec4 intJK = texelFetch(u_intStack,
                 ivec3(0,0,j*u_N+k), 0);
    vec4 intKI = texelFetch(u_intStack,
                 ivec3(0,0,k*u_N+i), 0);
    float sJK = intJK.g * cos(intJK.b);
    float sKI = intKI.g * cos(intKI.b);
    float quality = scoreIJ * sJK * sKI;
    float totPh = intIJ.b+intJK.b+intKI.b;
    // Resonance: Phi_L = 2*pi*m
    float valid = step(abs(mod(totPh,
                  2.0*PI)), PI*0.25);
    float cs = quality * valid;
    if (cs > bestQ) {
      bestQ = cs; bestPh = totPh;
    }
  }

  gl_FragColor = vec4(scoreIJ, bestPh,
                step(0.1,bestQ), bestQ);
}
\end{lstlisting}

\subsection{Pass 4: Readback and Metric Extraction}
\label{subsec:pass4}

\textbf{Purpose}: Validate via S-entropy conservation, extract final
match score, confidence, and circuit count. Transfer results from GPU
to CPU.

\textbf{Input}: $\mathcal{T}_{\text{int}}$, $\mathcal{T}_{\text{circ}}$.

\textbf{Output}: Validation texture $\mathcal{T}_{\text{valid}}$
(RGBA32F). R: match confidence. G: entropy residual. B: conservation
flag. A: validated match score.

\textbf{Shader pseudocode}:
\begin{lstlisting}[style=glsl]
// Pass 4: Readback / Validation
uniform sampler2D u_interference;
uniform sampler2D u_circuitMap;
uniform float u_epsilonEntropy;

void main() {
  vec2 uv = gl_FragCoord.xy / u_resolution;
  vec4 intD = texture(u_interference, uv);
  float ITot = intD.r;
  float V = intD.g;
  float dPhi = intD.b;

  // Local S-entropy
  float pOm = ITot;
  float Sk = -pOm * log(pOm + EPS);
  float pPh = V;
  float St = -pPh * log(pPh + EPS);
  float Se = 1.0 - Sk - St;
  float residual = abs(Sk + St + Se - 1.0);
  float conserved = step(residual,
                    u_epsilonEntropy);

  // Match with conservation gate
  float raw = V * cos(dPhi);
  float conf = raw * conserved;

  // Circuit boost
  vec4 circ = texture(u_circuitMap, uv);
  conf *= (1.0 + circ.b * circ.a);
  conf = clamp(conf, -1.0, 1.0);

  gl_FragColor = vec4(conf, residual,
                      conserved, conf);
}
\end{lstlisting}

The CPU-side readback uses \texttt{glReadPixels} (WebGL\,2),
\texttt{vkCmdCopyImageToBuffer} (Vulkan), or
\texttt{MTLBlitCommandEncoder} (Metal) to transfer the final texture
to CPU memory, from which the match score, confidence, and circuit
count are extracted as scalar aggregates.

% ── Figure 7 proposal ──────────────────────────────────────────────
% FIGURE 7: The five-pass pipeline. Flowchart:
% Raw Data -> [Pass 0: Encode] -> Spectral Image
% -> [Pass 1: Partition + Wavefunction] -> Psi texture
% -> [Pass 2: Interfere] -> Interference Pattern
% -> [Pass 3: Circuit] -> Circuit Map
% -> [Pass 4: Readback] -> {score, confidence, circuits}
% Each pass as a box with GPU icon, textures as arrows.

\begin{figure*}[!htbp]
    \centering
    \includegraphics[width=\textwidth]{figures/panel4_gpu_supervised_training.png}
    \caption{\textbf{GPU-Supervised Training: Convergence Without Human Labels.}
    (\textbf{A})~Three-dimensional surface plot of the composite loss landscape as a function of two probe parameters $w_1$ and $w_2$. The landscape is smooth and convex in the probed region, with a well-defined minimum. This smoothness is a prerequisite for gradient-based optimization and validates the Black-Box Oracle Convergence Theorem (Theorem~\ref{thm:convergence}): even though the GPU observation apparatus is a non-differentiable black box, the loss landscape induced by the physical observables is smooth enough for effective optimization.
    (\textbf{B})~Training loss curves over 100 iterations for two approaches: GPU-supervised training (blue, using only physical observables $P_{\mathrm{sharp}}$, $N_{\mathrm{obs}}$, $C_{\mathrm{phase}}$ as loss terms) and supervised baseline (orange, using ground-truth Dice as loss). Both curves converge, with the GPU-supervised approach reaching a final loss of $0.050$ and the supervised baseline reaching $0.009$. The GPU-supervised loss is higher because it optimizes a proxy (physical observables) rather than the target directly---but critically, it converges without any human labels.
    (\textbf{C})~Bar chart comparing final loss values: GPU-supervised ($0.050$) vs.\ supervised ($0.009$). The $5.8\times$ gap represents the cost of label-free training. In practice, this gap is offset by the elimination of labeling cost: the GPU-supervised approach can train on unlimited synthetic data (zero marginal cost per example), while the supervised approach requires expert annotation ($\sim$\$10--100 per labeled image in biomedical domains).
    (\textbf{D})~Observable improvement over training iterations: partition sharpness and phase coherence increase monotonically as the probe learns to generate better observation sequences, while noise decreases. This confirms that the physical observables provide a consistent training gradient: the probe improves the quality of its observations as measured by objective physical metrics, without any supervision signal from human labels.}
    \label{fig:panel_training}
\end{figure*}


% ═════════════════════════════════════════════════════════════════════
\section{Domain-Specific Encoders}
\label{sec:domain_encoders}

The spectral matching framework requires a domain encoder that maps raw
input data to the spectral image representation. This encoder (Pass~0)
is the only domain-specific component. We specify encoders for seven
domains.

\subsection{Microscopy Images}
\label{subsec:enc_micro}

\textbf{Input}: Image patch $I(x,y) \in \R^{W \times H}$.

\textbf{Encoding}: Compute 2D FFT \citep{cooley1965}; convert to polar
coordinates $(k, \theta)$; assign $n = \lceil k/\Delta k \rceil$,
$\ell = \lfloor \theta(n-1)/\pi \rfloor$; amplitude from magnitude
$|\hat{I}|$; construct spectral image.

\subsection{Molecular Structures (Vibrational Modes)}
\label{subsec:enc_mol}

\textbf{Input}: Absorption/emission spectrum $S(\nu)$ or computed
vibrational modes from a force-field model
\citep{lennardjones1924}.

\textbf{Encoding}: Identify peaks $\{\nu_k, I_k, \Gamma_k\}$; assign
$n = \lceil \nu_k / \Delta\nu \rceil$; angular partition from
vibrational symmetry; amplitude $A_k = \sqrt{I_k / \sum I_j}$.

\subsection{Chromatographic Peaks}
\label{subsec:enc_chrom}

\textbf{Input}: Chromatogram $C(t_R)$ (intensity vs.\ retention time).

\textbf{Encoding}: Detect peaks $\{t_{R,k}, A_k, w_k\}$; partition
depth $n_k = t_{R,k}^2 / (16 w_k^2)$ (theoretical plate model);
$\ell$ from peak asymmetry.

\subsection{Time Series / Signals}
\label{subsec:enc_ts}

\textbf{Input}: Time series $x(t)$ sampled at rate $f_s$.

\textbf{Encoding}: Compute STFT $X(f,t) = \int x(\tau)w(\tau-t)
e^{-2\pi i f\tau} d\tau$ \citep{oppenheim1999, harris1978}. The STFT
magnitude is already a 2D image (spectrogram). Assign partition
coordinates from frequency and time bins. This is the simplest case:
the spectrogram IS the spectral image (modulo normalisation). Compare
with dynamic time warping \citep{sakoe1978, berndt1994}, which the
interference metric subsumes.

\subsection{Text and Sequences}
\label{subsec:enc_text}

\textbf{Input}: Embedding vector $\mathbf{e} \in \R^D$ from a language
model \citep{vaswani2017, reimers2019}.

\textbf{Encoding}: Reshape to matrix
$\mathbf{E} \in \R^{\sqrt{D} \times \sqrt{D}}$; SVD
$\mathbf{E} = \mathbf{U\Sigma V}^T$ \citep{golub1965}; partition from
singular values ($n = k$, amplitude $A_k = \sigma_k / \|\Sigma\|_F$);
construct spectral image via learned projection.

\subsection{Genomic Sequences}
\label{subsec:enc_genomic}

\textbf{Input}: DNA/RNA sequence $s = (s_1, s_2, \ldots, s_L)$ over
alphabet $\{A, C, G, T\}$.

\textbf{Encoding}:
\begin{enumerate}[nosep]
  \item Compute $k$-mer frequency spectrum: for each $k$-mer of length
    $k$ ($k = 3, 4, 5, \ldots$), count occurrences normalised by
    sequence length.
  \item The $k$-mer frequency vector has dimension $4^k$. Reshape into
    a 2D array ($4^{k/2} \times 4^{k/2}$ for even $k$).
  \item Compute 2D FFT of the reshaped $k$-mer matrix.
  \item Assign partition coordinates from the Fourier modes: $n$ from
    spatial frequency, $\ell$ from angular character.
  \item Amplitude from Fourier magnitude.
\end{enumerate}
This provides a spectral fingerprint of the sequence that captures
both composition ($k=1$: nucleotide frequencies) and structure
($k \geq 3$: motif patterns). The spectral match between two sequences
correlates with alignment-based scores from BLAST \citep{altschul1990}
but is computed in $\BigO(D/P)$ rather than $\BigO(LD')$ where $L$ is
sequence length and $D'$ is database size.

\subsection{Graph Structures}
\label{subsec:enc_graph}

\textbf{Input}: Graph $G = (V, E)$ with adjacency matrix $\mathbf{A}$.

\textbf{Encoding}: Compute normalised Laplacian
$\mathbf{L} = \mathbf{I} - \mathbf{D}^{-1/2}\mathbf{A}\mathbf{D}^{-1/2}$
\citep{chung1997}; eigendecompose
$\mathbf{L} = \mathbf{U\Lambda U}^T$; partition $n = k$
(eigenvalue index), $\ell$ from nodal count of $\mathbf{u}_k$; amplitude
$A_k = \sqrt{\lambda_k / \sum_j \lambda_j}$. Connections to spectral
graph theory \citep{chung1997} and GNNs \citep{kipf2017} are direct.

% ── Figure 8 proposal ──────────────────────────────────────────────
% FIGURE 8: Domain encoder gallery. Seven sub-panels (microscopy,
% molecular, chromatography, time series, text, genomic, graph),
% each showing raw input -> encoder -> spectral image. All spectral
% images in the same coordinate system.


% ═════════════════════════════════════════════════════════════════════
\section{The Module API}
\label{sec:module_api}

\begin{definition}[Module API]
\label{def:api}
\begin{lstlisting}[style=api]
SpectralMatcher {

  // Encode input data as spectral image
  encode(data: Tensor, domain: DomainType)
    -> SpectralImage {
      amplitude: Texture2D<float>,
      phase: Texture2D<float>,
      partition: Texture2D<vec4>,
      wavefunction: Texture2D<vec2>
    }

  // Compare two spectral images
  match(img_a: SpectralImage,
        img_b: SpectralImage)
    -> MatchResult {
      score: float,        // [-1, 1]
      confidence: float,   // [0, 1]
      circuits: int,       // circuit count
      visibility: float,   // [0, 1]
      phaseDiff: float,    // [0, 2*pi)
      entropy: vec3,       // (Sk, St, Se)
      interferenceMap: Texture2D<vec4>
    }

  // Batch comparison with ranked results
  batch_match(query: SpectralImage,
              database: SpectralImage[])
    -> RankedResults {
      rankings: (int, float)[],
      graph: HarmonicGraph,
      circuits: Circuit[]
    }
}

DomainType = enum {
  MICROSCOPY, MOLECULAR_SPECTRA,
  CHROMATOGRAPHY, TIME_SERIES,
  TEXT_EMBEDDING, GENOMIC_SEQUENCE,
  GRAPH_STRUCTURE
}
\end{lstlisting}
\end{definition}

\begin{remark}[Usage Pattern]
The typical workflow is:
(1)~\texttt{encode} each input once (amortised cost);
(2)~\texttt{match} pairs in real-time ($\BigO(1)$ per pair);
(3)~\texttt{batch\_match} for collection-level ranked retrieval.
The encode step is performed once per datum; subsequent comparisons
reuse the cached spectral image (Corollary~\ref{cor:non_demolition}).
\end{remark}


% ═════════════════════════════════════════════════════════════════════
\section{Framebuffer Layout and Platform Notes}
\label{sec:framebuffer}

\begin{table}[t]
\centering
\caption{Texture specifications per pipeline pass.}
\label{tab:textures}
\begin{tabular}{lccc}
\toprule
\textbf{Pass} & \textbf{Name} & \textbf{Format} & \textbf{Size} \\
\midrule
0 & $\mathcal{T}_{\text{spec}}$ & RGBA32F & $N_\omega \!\times\! N_\varphi$ \\
1 & $\mathcal{T}_{\Psi}$ & RGBA32F & $N_\omega \!\times\! N_\varphi$ \\
2 & $\mathcal{T}_{\text{int}}$ & RGBA32F & $N_\omega \!\times\! N_\varphi$ \\
3 & $\mathcal{T}_{\text{circ}}$ & RGBA32F & $N \times N$ \\
4 & $\mathcal{T}_{\text{valid}}$ & RGBA32F & $N_\omega \!\times\! N_\varphi$ \\
\bottomrule
\end{tabular}
\end{table}

\textbf{WebGL\,2} \citep{webgl2spec}: Use
\texttt{WEBGL\_color\_buffer\_float} for RGBA32F render targets. Texture
arrays via \texttt{TEXTURE\_2D\_ARRAY} for Pass~3. Full-screen quad
rendering for all passes (compute shaders unavailable). Max texture
$4096 \times 4096$.

\textbf{Vulkan} \citep{vulkanspec}: Compute pipelines for Passes~0--1
and~4; graphics pipelines for Passes~2--3. Storage images with
\texttt{VK\_FORMAT\_R32G32B32A32\_SFLOAT}. Pipeline barriers between
passes.

\textbf{Metal}: Compute kernels for all passes. Threadgroup size
$16 \times 16$. Format \texttt{MTLPixelFormatRGBA32Float}.

\textbf{Memory budget} ($N_\omega = N_\varphi = 128$, $N = 1000$):
Per-item textures: $128 \times 128 \times 16 = 256$\,KB, two textures
per item = 512\,KB, total 512\,MB for 1000 items. Circuit map:
$1000 \times 1000 \times 16 = 16$\,MB. Well within commodity GPU memory
(8--24\,GB).


% ═════════════════════════════════════════════════════════════════════
%                    PART IV: VALIDATION AND APPLICATIONS
% ═════════════════════════════════════════════════════════════════════

\section{Application: Bioimage Matching}
\label{sec:app_bioimage}

\textbf{Dataset}: BBBC039 \citep{ljosa2012, caicedo2019}---fluorescence
microscopy of U2OS cell nuclei with diverse morphologies.

\textbf{Task}: Given a query nuclear patch, find all nuclei with matching
morphology.

\textbf{Protocol}: Encode all $\sim$200{,}000 patches via the microscopy
encoder ($n_{\max} = 32$, spectral image $32 \times 63$). Compute
interference with all database entries. Rank by match score.

\textbf{Baselines}: SIFT \citep{lowe2004}, ORB \citep{rublee2011},
DINOv2 \citep{oquab2023}.

\textbf{Metrics}: Precision@$k$ ($k = 1, 5, 10$), mAP, retrieval time.

\textbf{Expected outcomes}: Comparable precision to DINOv2; superior to
SIFT/ORB (which struggle with textureless nuclear morphologies);
orders-of-magnitude faster retrieval; physically interpretable
interference patterns.


% ═════════════════════════════════════════════════════════════════════
\section{Application: Molecular Spectral Matching}
\label{sec:app_molecular}

\textbf{Dataset}: NIST Chemistry WebBook \citep{nist2023}: $\sim$16{,}000
IR spectra, $\sim$10{,}000 Raman spectra.

\textbf{Task}: Identify unknown spectra by matching against the database.

\textbf{Protocol}: Encode all spectra via the molecular encoder. For each
query, compute interference. Rank by match score.

\textbf{Baselines}: Pearson correlation, Euclidean distance, Tanimoto
similarity on molecular fingerprints \citep{willett1998}.

\textbf{Expected outcomes}: Natural handling of peak shifts (Gaussian
kernel provides tolerance); partial matches reveal shared functional
groups; partition hierarchy respects fundamental vs.\ overtone structure.


% ═════════════════════════════════════════════════════════════════════
\section{Application: Signal Matching}
\label{sec:app_signal}

\textbf{Task}: Compare time series via spectral interference as an
alternative to dynamic time warping (DTW) \citep{sakoe1978}.

\textbf{Protocol}: Encode each time series via STFT
(Section~\ref{subsec:enc_ts}). The spectrogram IS the spectral image.
Compute interference.

\textbf{Baselines}: DTW \citep{sakoe1978, berndt1994}, which has
$\BigO(D^2)$ complexity per pair.

\textbf{Expected outcomes}: Spectral matching achieves $\BigO(D/P)$ per
pair vs.\ $\BigO(D^2)$ for DTW---a speedup of $D \cdot P$. For
$D = 1000$ and $P = 10^4$, this is $10^7 \times$ faster. The
interference metric captures frequency-domain similarity directly,
which DTW approximates through time warping.


% ═════════════════════════════════════════════════════════════════════
\section{Application: Cross-Domain Matching}
\label{sec:app_cross}

\textbf{Task}: Match a microscopy image of a fluorescently labelled
structure against a database of fluorophore emission spectra.

\textbf{Rationale}: Both the image and the emission spectrum are
represented as spectral images in the same $(\omega, \varphi)$ plane.
The fundamental frequencies of the fluorophore should produce
constructive interference at low $n$.

\textbf{Expected outcomes}: Noisier than within-domain matching, but the
shared spectral structure should be detectable. This capability is
qualitatively unique: no existing method can directly compare objects
from different measurement modalities.

% ── Figure 9 proposal ──────────────────────────────────────────────
% FIGURE 9: Cross-domain matching. Left: GFP microscopy image and
% its spectral image. Right: GFP emission spectrum and its spectral
% image. Centre: interference pattern (constructive at low freq).
% Bottom: mismatched fluorophore (destructive).

\begin{figure*}[!htbp]
    \centering
    \includegraphics[width=\textwidth]{figures/panel2_cross_domain_matching.png}
    \caption{\textbf{Cross-Domain Spectral Matching.}
    (\textbf{A})~Three-dimensional bar chart of mean match scores for all 16 domain pairs. The four domains are: molecular spectra (Mol), image patches (Img), time series (TS), and random sequences (Seq), each with 20 items. Blue bars along the diagonal represent within-domain matches; orange bars represent cross-domain matches. Within-domain means ($0.637$) systematically exceed cross-domain means ($0.534$), confirming that the spectral image representation preserves domain-specific structure while enabling universal comparison.
    (\textbf{B})~Full $80 \times 80$ match matrix (heatmap, RdYlBu colourmap) for all pairwise comparisons. Black grid lines delineate the four $20 \times 20$ domain blocks. The diagonal blocks (within-domain) are visibly warmer (higher match scores) than the off-diagonal blocks (cross-domain), confirming domain coherence. The matrix is perfectly symmetric by construction, and the diagonal entries equal~$1.0$ (self-match).
    (\textbf{C})~Box-and-whisker plots of match score distributions for all 10 unique domain pairs (4 within + 6 cross). Within-domain distributions (blue) have higher medians and tighter interquartile ranges than cross-domain distributions (orange). The separation between within- and cross-domain medians ($\Delta \approx 0.10$) provides the discriminative power for domain-aware retrieval.
    (\textbf{D})~Grouped bar chart comparing within-domain (blue) and cross-domain (orange) mean match scores per domain. All four domains show the same pattern: within $>$ cross. Molecular spectra show the largest separation ($\Delta = 0.13$), reflecting the sharp spectral peaks of vibrational frequencies. Time series show the smallest separation ($\Delta = 0.07$), reflecting the broader spectral content of temporal signals.}
    \label{fig:panel_cross_domain}
\end{figure*}


% ═════════════════════════════════════════════════════════════════════
\section{Scalability Analysis}
\label{sec:scalability}

\begin{theorem}[Throughput]
\label{thm:throughput}
On a modern GPU with $P$ cores operating at clock frequency $f_c$ and
memory bandwidth $B$, the spectral matching throughput for batch
pairwise comparison is:
\begin{equation}
\label{eq:throughput}
  \Theta = \min\left(\frac{P \cdot f_c}{5 \cdot C_{\text{pass}}},
    \frac{B}{2 \cdot D \cdot 4}\right) \text{ comparisons/s},
\end{equation}
where $C_{\text{pass}}$ is the average instruction count per pixel per
pass (typically 20--50 instructions), the factor 5 accounts for five
pipeline passes, $D$ is the spectral image size in pixels, and the
factor $2 \cdot D \cdot 4$ accounts for reading two RGBA32F textures.
\end{theorem}

\begin{proof}
The compute-bound limit is $P$ cores each executing
$f_c / (5 \cdot C_{\text{pass}})$ full pipeline evaluations per second
(5 passes, $C_{\text{pass}}$ cycles each). The memory-bound limit is
$B / (2 \cdot D \cdot 4 \cdot 4)$ where $4 \cdot 4 = 16$ bytes per
RGBA32F texel and two textures are read per comparison. The throughput
is the minimum (bottleneck).
\end{proof}

\begin{example}[RTX 4090 Throughput]
NVIDIA RTX 4090: $P = 16384$ cores, $f_c = 2.52$\,GHz,
$B = 1008$\,GB/s. With $D = 128 \times 128 = 16384$ pixels,
$C_{\text{pass}} = 30$ instructions:
\begin{align}
  \Theta_{\text{compute}} &= \frac{16384 \times 2.52 \times 10^9}{5
    \times 30 \times 16384} \approx 1.68 \times 10^7 \text{ comp/s}, \\
  \Theta_{\text{memory}} &= \frac{1008 \times 10^9}{2 \times 16384
    \times 16} \approx 1.93 \times 10^6 \text{ comp/s}.
\end{align}
The system is memory-bound at approximately $\mathbf{2 \times 10^6}$
\textbf{pairwise comparisons per second}. For $N = 1000$ items, all
$\binom{1000}{2} \approx 500{,}000$ pairwise comparisons complete in
$\sim$0.25\,s.
\end{example}

\begin{table}[t]
\centering
\caption{Predicted throughput on representative GPUs.}
\label{tab:throughput}
\begin{tabular}{lrrr}
\toprule
\textbf{GPU} & \textbf{Cores} & \textbf{BW (GB/s)} & \textbf{Comp/s} \\
\midrule
RTX 3060 & 3584 & 360 & $6.9 \times 10^5$ \\
RTX 4090 & 16384 & 1008 & $1.9 \times 10^6$ \\
A100 & 6912 & 2039 & $3.9 \times 10^6$ \\
H100 & 14592 & 3350 & $6.4 \times 10^6$ \\
\bottomrule
\end{tabular}
\end{table}


% ═════════════════════════════════════════════════════════════════════
\section{Validation Framework}
\label{sec:validation}

\subsection{Internal Consistency}

\begin{enumerate}[nosep]
  \item \textbf{S-entropy conservation}: $|\Sk + \St + \Se - 1| < 10^{-4}$
    for every comparison. Failure rate $< 0.01\%$.
  \item \textbf{Self-matching}: $\Match(A, A) = 1.0 \pm 10^{-6}$.
  \item \textbf{Symmetry}: $|\Match(A,B) - \Match(B,A)| < 10^{-6}$.
  \item \textbf{Circuit closure}: $|\sum_i \Delta\varphi_i - 2\pi m|
    < 0.1$ for all detected circuits.
  \item \textbf{Triangle inequality}: Verify
    \eqref{eq:circuit_inequality} for all triples.
\end{enumerate}

\subsection{External Metrics}

\begin{table}[t]
\centering
\caption{Target validation metrics per application.}
\label{tab:validation}
\begin{tabular}{lll}
\toprule
\textbf{Application} & \textbf{Metric} & \textbf{Target} \\
\midrule
Bioimage & Precision@10 & $\geq 0.85$ \\
         & mAP & $\geq 0.80$ \\
         & 1000 queries & $< 1$\,s \\
\midrule
Molecular & Top-1 accuracy & $\geq 0.90$ \\
          & Top-5 accuracy & $\geq 0.95$ \\
\midrule
Signal & Spearman $\rho$ vs DTW & $\geq 0.90$ \\
       & Speedup vs DTW & $\geq 10^3\times$ \\
\midrule
Cross-domain & Top-3 accuracy & $\geq 0.70$ \\
\bottomrule
\end{tabular}
\end{table}

\subsection{Ablation Studies}

\begin{enumerate}[nosep]
  \item Vary $n_{\max}$ from 8 to 256: quality vs.\ cost.
  \item Compare Gaussian, Hann, Blackman windows \citep{harris1978}:
    selectivity vs.\ leakage.
  \item Full complex vs.\ magnitude-only: quantify phase information
    value.
  \item Circuit length 3 to 7: precision gain from circuit validation.
\end{enumerate}

% ── Figure 10 proposal ────────────────────────────────────────────
% FIGURE 10: Ablation results (proposed). (a) Match quality vs
% n_max showing saturation. (b) Window function comparison.
% (c) Phase vs magnitude-only. (d) Circuit length vs precision.

\begin{figure*}[!htbp]
    \centering
    \includegraphics[width=\textwidth]{figures/panel4_s_entropy_conservation.png}
    \caption{\textbf{S-Entropy Conservation in Spectral Representations.}
    (\textbf{A})~Three-dimensional scatter plot of S-entropy coordinates $(S_k, S_t, S_e)$ for all 50 synthetic systems, with the conservation manifold $S_k + S_t + S_e = 1$ rendered as a semi-transparent surface. All data points lie exactly on the manifold, confirming S-entropy conservation to machine precision (maximum deviation $1.11 \times 10^{-16}$). The cluster position reflects the spectral structure of the synthetic systems: $S_k$ (kinetic/knowledge entropy) encodes spectral diversity, $S_t$ (temporal entropy) encodes frequency centroid, and $S_e$ (evolution entropy) encodes spectral concentration.
    (\textbf{B})~Stacked bar chart of entropy partitioning per system. Blue bars represent $S_k$, orange $S_t$, green $S_e$. Every bar sums to exactly $1.0$ (black dashed line), visually confirming pixel-level conservation across all 50 systems. The variation in proportions across systems reflects the diversity of synthetic frequency configurations---systems with many dispersed components have higher $S_k$; systems with concentrated energy at few frequencies have higher $S_e$.
    (\textbf{C})~Histogram of conservation deviations $|S_k + S_t + S_e - 1|$. All 50 deviations are numerically zero at IEEE~754 double-precision limits, producing a single-bin distribution at $0.0$. This confirms that the conservation constraint is enforced exactly by construction, not approximately by optimization---a direct consequence of defining $S_e = 1 - S_k - S_t$ as the closure residual.
    (\textbf{D})~Donut chart of mean entropy partitioning across all systems. The proportions quantify where spectral information resides on average, providing a compact summary of the ensemble's categorical state distribution. These proportions are system-dependent: real molecular spectra would show different partitioning from genomic k-mer spectra, but the conservation law $S_k + S_t + S_e = 1$ holds universally across all domains.}
    \label{fig:panel_entropy}
\end{figure*}

% ═════════════════════════════════════════════════════════════════════
\section{Discussion}
\label{sec:discussion}

\subsection{The Computer Vision Unification}

The spectral matching framework reveals that computer vision is not
merely a subfield of AI that processes images. It is the
\emph{fundamental computational paradigm} for comparing any systems that
can be represented as oscillatory processes in bounded phase space---which,
by Theorem~\ref{thm:oscillatory_necessity}, includes all persistent
observable systems. Images are not special inputs; they are the universal
representation.

\subsection{Plugging In New Encoders}

The architecture is deliberately modular. To add support for a new data
domain, a collaborator need only:
\begin{enumerate}[nosep]
  \item Define the mapping from raw data to $(\omega_k, A_k, \phi_k)$
    triples (the domain encoder).
  \item Implement Pass~0 for the new domain type.
  \item Validate via the internal consistency checks
    (Section~\ref{sec:validation}).
\end{enumerate}
Passes 1--4 are domain-independent and require no modification. The
module API (Section~\ref{sec:module_api}) accepts the new domain type as
an enum value. This plug-in architecture means the framework scales to
arbitrary new domains without re-engineering the comparison engine.

\subsection{Unification of Existing Methods}

The spectral matching module replaces, in principle, the following
domain-specific matching methods with a single unified framework:
\begin{itemize}[nosep]
  \item SIFT \citep{lowe2004} and ORB \citep{rublee2011} for image
    matching.
  \item DTW \citep{sakoe1978} for time series matching.
  \item BLAST \citep{altschul1990} for sequence alignment.
  \item Tanimoto/fingerprint matching \citep{willett1998} for molecular
    similarity.
  \item Cosine similarity for embedding comparison.
  \item Graph kernel methods \citep{xu2019} for graph comparison.
\end{itemize}
Each of these is a special case of spectral interference, restricted to
a particular domain encoder and comparison metric. The spectral framework
subsumes them all (Theorem~\ref{thm:metric_equivalence}).

\subsection{The Spectral Microscope Metaphor}

Every comparison performed by this module can be visualised as looking
at two objects through a ``spectral microscope'': the encoder converts
each object to a pattern of light (the spectral image), and the GPU
interferometer superimposes the two patterns. Where the patterns align,
one sees bright fringes (match). Where they differ, one sees dark bands
(mismatch). The interference pattern is not merely a scalar score but a
spatially resolved map of agreement and disagreement across the spectral
domain. This is fundamentally more informative than any single-number
similarity metric.

\subsection{Neural Networks as Degenerate Interference}

A ReLU neural network computes
$\mathbf{y} = \max(0, \mathbf{W}\mathbf{x} + \mathbf{b})$---a
phase-insensitive superposition. This is a degenerate case of
interference where all phases are zero. The spectral framework uses full
complex-valued interference. This predicts that complex-valued networks
\citep{trabelsi2018} should outperform real-valued networks for
comparison tasks, and that attention mechanisms \citep{vaswani2017}
(which compute degenerate spectral matching) could be enhanced by
incorporating explicit phase information.

\subsection{Relationship to Kuramoto Synchronisation}

The Kuramoto model \citep{kuramoto1975, strogatz2000} describes coupled
oscillator synchronisation. Matching circuit detection (Pass~3) identifies
synchronised clusters: spectral images with locked phases
($\Delta\varphi \to 0$) form matching circuits, analogous to phase-locked
oscillator populations above the Kuramoto coupling threshold.

\subsection{Limitations}

\textbf{Unbounded systems.} Axiom~\ref{ax:bounded} excludes genuinely
unbounded systems. In practice, all computationally representable
systems are bounded.

\textbf{Encoder quality.} The framework's quality depends on the
domain-specific encoder. A poor encoder yields uninformative spectral
images. The theoretical guarantee is for a perfect encoder; practical
encoders approximate this ideal.

\textbf{GPU memory.} Batch comparison of $N$ items requires storing
$N$ spectral images. For $N > 10^6$ with high-resolution images,
multi-GPU or streaming approaches are necessary.

\textbf{All-pairs scaling.} While single-pair comparison is $\BigO(1)$,
all-pairs remains $\BigO(N^2)$. For very large $N$, approximate
methods (LSH in spectral space, hierarchical harmonic grouping) are
needed.

\subsection{Future Directions}

\textbf{Learned spectral encoders.} Train neural networks to learn
optimal domain encoders end-to-end, using the conservation loss
(Proposition~\ref{prop:conservation_loss}) as a regulariser.

\textbf{Real-time streaming.} Implement the pipeline in a streaming
mode where new data are encoded and compared against a rolling window
of spectral images, enabling real-time anomaly detection.

\textbf{Quantum implementation.} On a quantum computer, each spectral
component is a genuine quantum amplitude, and interference is physical
rather than simulated, enabling comparison of exponentially large spectra.

\textbf{Custom silicon.} An ASIC implementing the five-pass pipeline in
fixed hardware would achieve sub-microsecond latency per comparison.

\textbf{Biological interpretation.} Neural oscillations at multiple
frequency bands \citep{buzsaki2006} may implement biological spectral
matching: sensory oscillations interfere with memory oscillations, and
constructive maxima drive recognition.


% ═════════════════════════════════════════════════════════════════════
\section{Conclusion}
\label{sec:conclusion}

We have established from two axioms a complete theoretical framework
and engineering blueprint proving that \emph{all comparison is computer
vision}. The argument chain is:

\begin{enumerate}[nosep]
  \item \textbf{Bounded phase space} (Axiom~\ref{ax:bounded}) and
    \textbf{categorical observation} (Axiom~\ref{ax:categorical})
    define the arena.
  \item \textbf{Oscillatory necessity}
    (Theorem~\ref{thm:oscillatory_necessity}): bounded + persistent
    + non-degenerate $\Rightarrow$ oscillatory.
  \item \textbf{Koopman spectral decomposition}
    (Theorem~\ref{thm:koopman_spectral}): every bounded oscillatory
    system has spectrum
    $\Spec(\mathcal{S}) = \{(\omega_k, A_k, \phi_k)\}$.
  \item \textbf{Partition coordinates}
    (Theorem~\ref{thm:partition_coords}): $(n, \ell, m, s)$ with
    $C(n) = 2n^2$.
  \item \textbf{Triple equivalence} (Theorem~\ref{thm:triple}):
    oscillation $\equiv$ enumeration $\equiv$ partition function,
    $S = \kb M \ln n$.
  \item \textbf{Fundamental identity} (Theorem~\ref{thm:fundamental}):
    $dM/dt = \omega/(2\pi) = 1/\langle\taup\rangle$.
  \item \textbf{Spectral Image Theorem} (Theorem~\ref{thm:spectral_image}):
    $\Spec \cong$ image, injective, metric-preserving.
  \item \textbf{Universal Reduction Theorem}
    (Theorem~\ref{thm:universal_reduction}):
    $d(X,Y) = d_{\mathrm{CV}}(I_X, I_Y)$ exactly.
  \item \textbf{S-entropy conservation}
    (Theorem~\ref{thm:s_conservation}): $\Sk + \St + \Se = 1$.
  \item \textbf{Commutation} (Theorem~\ref{thm:commutation}):
    $[\Ocat, \Ophys] = 0$.
  \item \textbf{Interference matching}
    (Theorem~\ref{thm:interference}): $V \to 1$ = match,
    $V \to 0$ = mismatch.
  \item \textbf{GPU = interference engine}
    (Theorem~\ref{thm:gpu_interference}): fragment shader $\cong$
    wave field evaluator.
  \item \textbf{Complexity} (Theorem~\ref{thm:complexity}):
    $\BigO(1)$ per comparison when $D \leq P$.
\end{enumerate}

The Universal Spectral Matching module converts any comparison
problem---regardless of domain---into GPU render passes that compute
wave interference between spectral images. It is domain-agnostic (via
pluggable encoders), hardware-efficient ($\sim 10^6$ comparisons/s on
commodity GPUs), theoretically exact (not an approximation), and
physically interpretable (interference patterns show where and how
objects match).

The deepest implication is this: because all bounded systems are
oscillatory, and oscillations have spectra, and spectra are images,
the comparison of \emph{anything} is, at bottom, a problem of
\emph{seeing}. Every comparison is a comparison of pictures. Every
matcher is a vision system. And every GPU is, in its fundamental
operation, a machine for creating and reading interference patterns.

\vspace{1em}
\noindent\textbf{Data and Code Availability.}
All theoretical results are self-contained in this paper. Shader source
code, API implementations, and benchmark scripts will be released at
the module repository upon publication.

\vspace{1em}
\noindent\textbf{Acknowledgments.}
The author thanks the open-source GPU computing community and the
developers of WebGL, Vulkan, and Metal for creating the substrates
on which spectral matching is realised.


% ═════════════════════════════════════════════════════════════════════
%                         APPENDICES
% ═════════════════════════════════════════════════════════════════════

\appendix

\section{Proof of Parseval's Theorem for Partition Coordinates}
\label{app:parseval}

\begin{theorem}
For a system with spectral decomposition
$\{(n_k, \ell_k, m_k, s_k, A_k)\}$ and categorical occupation numbers
$\{N_j\}$,
\begin{equation}
\label{eq:parseval}
  \sum_k |A_k|^2 = \frac{1}{M} \sum_j N_j = 1.
\end{equation}
\end{theorem}

\begin{proof}
By Parseval's theorem for Fourier series \citep{fourier1822},
$\sum_k |A_k|^2 = (1/T) \int_0^T |c(t)|^2 \, dt$ where $c(t)$ is the
categorical time series. Since amplitudes are normalised by $\sqrt{M}$,
the sum is $\sum_k |A_k|^2 = 1$. Categorically:
$|A_k|^2 = p_k$ and $\sum_k p_k = 1$.
\end{proof}


\section{Visibility--Mutual Information Derivation}
\label{app:visibility_mi}

\begin{theorem}
For spectral images with Gaussian-distributed amplitudes,
\begin{equation}
\label{eq:vis_mi}
  I(I_A; I_B) = -\frac{1}{2}\ln(1 - V_{\text{global}}^2).
\end{equation}
\end{theorem}

\begin{proof}
For jointly Gaussian random variables with correlation $\rho$,
$I(X;Y) = -\frac{1}{2}\ln(1 - \rho^2)$ \citep{cover2006}. The degree
of coherence $|\gamma_{12}|$ between two wave fields equals $V$ for
quasi-monochromatic fields \citep{born1999, vanzernike1938}. Identifying
$\rho = V_{\text{global}}$ yields the result.
\end{proof}


\section{Shader Implementation Notes}
\label{app:shader}

\textbf{GLSL version}: Shaders target GLSL ES 3.0 (WebGL\,2) or
GLSL 4.50 (Vulkan) \citep{glslspec}. Constants: \texttt{EPS} $= 10^{-7}$;
\texttt{SIGMA\_W} $= \Delta\omega / 2 = \pi / n_{\max}$;
\texttt{PI} $= 3.14159265$; \texttt{INV\_SQRT2} $= 0.70710678$.

\textbf{Precision}: 32-bit float throughout. For high-resolution
molecular spectra, 64-bit is available on Vulkan/CUDA.

\textbf{Synchronisation}: Each pass must complete before the next
(output becomes input). WebGL\,2 guarantees this by single-threaded draw
calls. Vulkan requires \texttt{vkCmdPipelineBarrier}. Metal uses
\texttt{MTLCommandBuffer} ordering.


\section{Summary of Theorems}
\label{app:theorem_summary}

\begin{enumerate}[nosep]
  \item \textbf{Oscillatory Necessity} (Thm.~\ref{thm:oscillatory_necessity}):
    Bounded + persistent + non-degenerate $\Rightarrow$ oscillatory.
  \item \textbf{Spectral Existence} (Cor.~\ref{cor:spectral_existence}):
    Oscillatory $\Rightarrow$ Fourier decomposition exists.
  \item \textbf{Partition Coordinates} (Thm.~\ref{thm:partition_coords}):
    $(n, \ell, m, s)$ with $C(n) = 2n^2$.
  \item \textbf{Capacity} (Thm.~\ref{thm:capacity}): $C(n) = 2n^2$.
  \item \textbf{Triple Equivalence} (Thm.~\ref{thm:triple}):
    Oscillation $\equiv$ enumeration $\equiv$ partition function.
  \item \textbf{Fundamental Identity} (Thm.~\ref{thm:fundamental}):
    $dM/dt = \omega/(2\pi) = 1/\langle\taup\rangle$.
  \item \textbf{Koopman Spectral} (Thm.~\ref{thm:koopman_spectral}):
    Finite spectrum via Koopman operator.
  \item \textbf{Spectral Image} (Thm.~\ref{thm:spectral_image}):
    $\Spec \cong$ image, injective, metric-preserving.
  \item \textbf{Universal Reduction} (Thm.~\ref{thm:universal_reduction}):
    $d_{\mathrm{part}}(X,Y) = d_{\mathrm{CV}}(I_X, I_Y)$.
  \item \textbf{S-Entropy Conservation} (Thm.~\ref{thm:s_conservation}):
    $\Sk + \St + \Se = 1$.
  \item \textbf{Commutation} (Thm.~\ref{thm:commutation}):
    $[\Ocat, \Ophys] = 0$.
  \item \textbf{Interference} (Thm.~\ref{thm:interference}):
    Two-beam interference formula.
  \item \textbf{Metric Equivalence} (Thm.~\ref{thm:metric_equivalence}):
    Match $\supseteq$ cosine, $L^2$, MI.
  \item \textbf{Circuit Consistency} (Thm.~\ref{thm:circuit_consistency}):
    Triangle inequality.
  \item \textbf{GPU--Interference} (Thm.~\ref{thm:gpu_interference}):
    Fragment shader $\cong$ interference engine.
  \item \textbf{Complexity} (Thm.~\ref{thm:complexity}):
    $\BigO(1)$ per comparison when $D \leq P$.
  \item \textbf{Throughput} (Thm.~\ref{thm:throughput}):
    $\sim 10^6$ comparisons/s on commodity GPU.
\end{enumerate}

% ═════════════════════════════════════════════════════════════════════

\bibliographystyle{plainnat}
\bibliography{references}

\end{document}
